\documentclass[12pt,a4paper,oneside]{article}
\usepackage{fontspec}
\usepackage{geometry}
\usepackage{xcolor}
\usepackage{tikz}
\usepackage{graphicx}
\usepackage{svg}
\usepackage{pagecolor}
\usepackage{setspace}
\usepackage{titlesec}
\usepackage{fancyhdr}
\usepackage{microtype}
\usepackage{amssymb}
\usepackage{tocloft}

% Suppress lost character log warnings globally
\tracinglostchars=0

% Page Geometry
\geometry{
  a4paper,
  left=25mm,
  right=25mm,
  top=28mm,
  bottom=28mm
}

% Colors - Matmaila / Vintage Parchment Theme
\definecolor{matmaila}{HTML}{C5B99D}      % Authentic matmaila earthy beige parchment
\definecolor{matmailadark}{HTML}{2B251F}  % Deep charcoal brown text
\definecolor{bordercolor}{HTML}{3D3428}   % Vintage dark brown frame

% Fonts Configuration
\setmainfont{Noto Serif Devanagari}[
  Renderer=HarfBuzz,
  Script=Devanagari,
  Language=Hindi,
  UprightFont=*-Regular,
  BoldFont=*-Bold
]

\newfontfamily\engfont{Latin Modern Roman}

% Configure itemize bullet to use engfont textbullet
\renewcommand{\labelitemi}{\engfont\textbullet}

% Universal TOC Page Number Box Width (4.5em reserved for Roman numerals like XXVII)
\makeatletter
\renewcommand{\@pnumwidth}{4.5em}
\renewcommand{\@tocrmarg}{5.0em}
\makeatother

\renewcommand{\cftsecpagefont}{\engfont}
\renewcommand{\cftsubsecpagefont}{\engfont}
\renewcommand{\cftsecfont}{\normalfont}
\renewcommand{\cftsubsecfont}{\normalfont}

% Line Spacing
\setstretch{1.35}

% Paragraph Formatting
\setlength{\parindent}{1.5em}
\setlength{\parskip}{0.5em}

% Section Titles Styling
\titleformat{\section}[block]
  {\centering\Large\bfseries\color{bordercolor}}
  {}{0pt}{}
  [\vspace{0.4em}{\color{bordercolor}\rule{0.35\textwidth}{0.8pt}}\vspace{0.8em}]

\titleformat{\subsection}[block]
  {\large\bfseries\color{bordercolor}}
  {}{0pt}{}

% Header and Footer - Roman Page Numbers
\pagestyle{fancy}
\fancyhf{}
\fancyfoot[C]{\color{bordercolor}\small{\engfont \thepage}}
\renewcommand{\headrulewidth}{0pt}
\renewcommand{\footrulewidth}{0pt}

% Redefine plain page style for TOC and chapter pages so footer uses engfont
\fancypagestyle{plain}{%
  \fancyhf{}%
  \fancyfoot[C]{\color{bordercolor}\small{\engfont \thepage}}%
  \renewcommand{\headrulewidth}{0pt}%
  \renewcommand{\footrulewidth}{0pt}%
}

\begin{document}
\pagenumbering{Roman}
\pagecolor{matmaila}
\color{matmailadark}

% ================= PAGE I: FRONT COVER =================
\begin{titlepage}
\thispagestyle{empty}

\begin{tikzpicture}[remember picture, overlay]
  % Outer Double Frame Lines
  \draw[line width=1.2pt, color=bordercolor] 
    ([yshift=-10mm, xshift=10mm]current page.north west) rectangle 
    ([yshift=10mm, xshift=-10mm]current page.south east);

  \draw[line width=0.5pt, color=bordercolor] 
    ([yshift=-12mm, xshift=12mm]current page.north west) rectangle 
    ([yshift=12mm, xshift=-12mm]current page.south east);

  % Top Decorative Ornamental Dots Bar
  \draw[line width=1.2pt, color=bordercolor] 
    ([yshift=-20mm, xshift=18mm]current page.north west) -- 
    ([yshift=-20mm, xshift=-18mm]current page.north east);
  
  \foreach \x in {-7.5,-6,-4.5,-3,-1.5,0,1.5,3,4.5,6,7.5} {
    \node[color=bordercolor, font=\engfont] at ([yshift=-17mm, xshift=\x cm]current page.north) {\small $\therefore$};
  }

  % Bottom Decorative Ornamental Dots Bar
  \draw[line width=1.2pt, color=bordercolor] 
    ([yshift=20mm, xshift=18mm]current page.south west) -- 
    ([yshift=20mm, xshift=-18mm]current page.south east);
  
  \foreach \x in {-7.5,-6,-4.5,-3,-1.5,0,1.5,3,4.5,6,7.5} {
    \node[color=bordercolor, font=\engfont] at ([yshift=23mm, xshift=\x cm]current page.south) {\small $\because$};
  }
\end{tikzpicture}

\vspace*{1.2cm}

\begin{center}
  % Main Header / Title Block
  {\Huge \bfseries \color{matmailadark} कम्युनिस्ट पार्टी के घोषणापत्र} \\[0.7em]
  {\large \color{bordercolor} {\engfont (The Communist Manifesto)}} \\[0.4em]
  {\small \color{bordercolor} भोजपुरी अनुवाद} \\[1.8cm]

  % Center Emblem: Direct SVG Integration using \includesvg
  \includesvg[width=6.0cm]{communist_emblem.svg}

  \vfill

  % Bottom Left Metadata Block
  \begin{flushleft}
  \hspace*{1.0cm}
  \begin{tikzpicture}
    \draw[line width=1.8pt, color=bordercolor] (0,0) -- (0, 2.4);
    \node[anchor=west, text width=12.5cm, color=matmailadark, font=\normalfont] at (0.35, 1.2) {
      \small
      {\bfseries मूल लेखक:} कार्ल मार्क्स आ फ्रेडरिक एंगेल्स \\
      {\bfseries भोजपुरी अनुवाद:} आदित्य कुमार \\
      {\bfseries भाषा शैली:} भोजपुरी \\
      {\bfseries परियोजना:} भोजपुरी रिवाइव प्रोजेक्ट {\engfont (Bhojpuri Revive Project)} \\
      {\bfseries संस्करण:} प्रथम भोजपुरी संस्करण (2026)
    };
  \end{tikzpicture}
  \end{flushleft}
\end{center}

\vspace*{0.8cm}
\end{titlepage}

% ================= PAGE II: DEDICATED COPYRIGHT PAGE =================
\newpage
\thispagestyle{empty}

\vfill
\begin{center}
  {\Large \bfseries \color{bordercolor} कम्युनिस्ट पार्टी के घोषणापत्र} \\[0.3em]
  {\small \color{bordercolor} {\engfont (The Communist Manifesto --- Bhojpuri Translation)}} \\[1.5em]
  {\color{bordercolor}\rule{0.4\textwidth}{0.8pt}}
\end{center}

\vspace*{1.2cm}

\begin{center}
\begin{minipage}{0.88\textwidth}
  \small
  \setstretch{1.3}
  \setlength{\parskip}{0.7em}

  {\bfseries पुस्तक का नाम:} कम्युनिस्ट पार्टी के घोषणापत्र \\
  {\bfseries मूल रचयिता:} कार्ल मार्क्स आ फ्रेडरिक एंगेल्स (1848) \\
  {\bfseries भोजपुरी अनुवादक:} आदित्य कुमार {\engfont (Aditya Kumar)} \\
  {\bfseries भाषा शैली:} भोजपुरी \\
  {\bfseries सर्वाधिकार:} {\engfont \copyright} 2026 आदित्य कुमार (सर्वहक़ सुरक्षित) \\
  {\bfseries प्रथम संस्करण:} 2026

  \vspace*{1.0em}
  {\color{bordercolor}\rule{\textwidth}{0.4pt}}
  \vspace*{1.0em}

  {\bfseries सर्वाधिकार आ लाइसेंस संबंधी शर्त {\engfont (Copyright Notice)}:}

  \begin{itemize}\setlength{\itemsep}{0.4em}
    \item {\bfseries व्यक्तिगत उपयोग {\engfont (Personal Use)}:} ई भोजपुरी अनुवाद केवल व्यक्तिगत पठन, निजी अध्ययन आ गैर-व्यावसायिक पठन खातिर मुफ़्त उपलब्ध बा।
    \item {\bfseries व्यावसायिक आ अकादमिक छपाई {\engfont (Commercial \& Academic Mass Printing)}:} कवनो प्रकार के व्यावसायिक प्रकाशन, बड़े पैमाने पर छपाई {\engfont (Mass Printing)}, व्यावसायिक संस्थान या अकादमिक संस्थाओं द्वारा बड़े पैमाने पर वितरण या व्यावसायिक उपयोग खातिर अनुवादक (आदित्य कुमार) से लिखित पूर्वानुमति आ रॉयल्टी {\engfont (Royalty Agreement)} लेहल अनिवार्य बा।
  \end{itemize}

  \vspace*{0.8em}
  {\color{bordercolor}\rule{\textwidth}{0.4pt}}
  \vspace*{0.8em}

  \begin{center}
    {\engfont \copyright{} 2026 Aditya Kumar. All Rights Reserved.} \\
    {\small {\engfont Free for personal non-commercial study only. Commercial printing, mass distribution, or commercial academic use requires prior written permission and royalty licensing.}} \\
    \vspace*{0.5em}
    {\small स्वतंत्र अनुवाद परियोजना | भोजपुरी रिवाइव प्रोजेक्ट {\engfont (Bhojpuri Revive Project)}}
  \end{center}
\end{minipage}
\end{center}
\vfill

% ================= PAGE III: TABLE OF CONTENTS (विषय-सूची - SINGLE PAGE) =================
\newpage
\thispagestyle{empty}

\begin{center}
  {\Large \bfseries \color{bordercolor} विषय-सूची} \\[0.4em]
  {\color{bordercolor}\rule{0.3\textwidth}{0.8pt}}
\end{center}

\vspace*{0.6cm}

\renewcommand{\contentsname}{}
\setcounter{tocdepth}{2}
{\makeatletter
\renewcommand{\@dotsep}{4.5}
\makeatother
\tableofcontents}

\vfill
\begin{center}
  {\small \color{bordercolor} भोजपुरी रिवाइव प्रोजेक्ट | प्रथम भोजपुरी संस्करण (2026)}
\end{center}

\newpage

% ================= MAIN BODY CONTENT =================
\cleardoublepage
\cleardoublepage
\section*{जीवनी}
\addcontentsline{toc}{section}{जीवनी}
कार्ल हेनरिच मार्क्स के जनम 5 मई 1818 के जर्मनी के ट्रायर शहर में एगो खुशहाल मध्यमवर्गीय यहूदी परिवार में भइल रहे। उनकर पिता, जे एगो नामी वकील आ प्रबुद्ध उग्रवादी सोच के पक्षधर रहलन, मार्क्स के लइकाए में परिवार के लूथरन ईसाई धर्म में बदल दिहलन, ताकि प्रूशिया में ओह घरी यहूदियन पर होखे वाला भेदभाव आ प्रताड़ना से परिवार बच सके। मार्क्स के अपना पिता के देखरेख में एगो व्यापक आ उदार शिक्षा मिलल। आगे चल के उनकरा फ्रीहेर लुडविग वॉन वेस्टफ़ेलेन जइसन कुलीन प्रूशियन विचारक के रूप में एगो बौद्धिक मार्गदर्शक (गुरु) मिललन, जिनकरा संगे मार्क्स ओह ज़माना के बड़का-बड़का साहित्यिक आ दार्शनिक हस्तियन पर गहिर चर्चा करत रहलन। वेस्टफ़ेलेन हीं रहलन जे नौजवान मार्क्स के शुरुआती फ्रांसीसी समाजवादी विचारक सेंट-सिमोन के विचारन से परिचित करवले रहलन।

बॉन आ बर्लिन विश्वविद्यालय में पढ़ाई के दौरान मार्क्स हेगेल के दर्शन से बहुत गहिर रूप से प्रभावित भइलन। हालाँकि मार्क्स अपना हेगेलवादी प्रोफेसरन से काफी सीखलन, बाकिर अंत में ऊ "नौजवान हेगेलवादी" (यंग हेगेलियंस) नाम के छात्र-विचारकन के गुट के ओर खिंचा गइलन। डेविड स्ट्रॉस, ब्रूनो बॉअर आ मैक्स स्टिरनर जइसन क्रांतिकारी आ विद्रोही सोच वाला एह युवा लोगन के समूह हेगेल के विचारन से प्रेरणा त लेत रहे, बाकिर ऊ लोग हेगेल के दर्शन के सबसे उग्र आ क्रांतिकारी पक्ष के आगे बढ़ावे खातिर अडिग रहे। खास करके, ई वामपंथी हेगेलवादी लोग हेगेल के राजनीतिक आ धार्मिक दर्शन में छिपल रूढ़िवाद पर तीखा सवाल उठावत रहे। ओह समय मार्क्स के मन अकादमिक (विश्वविद्यालयी) प्रोफेसर बने के रहे, बाकिर उनकर उग्र राजनीतिक विचारन के कारण सरकारी नियंत्रण वाला विश्वविद्यालय प्रणाली में उनकरा खातिर नौकरी के दुआर बंद हो गइल। एह से मजबूर हो के मार्क्स पत्रकारिता के राह पकड़लन, जहाँ उनकर तीखा आ क्रांतिकारी लेखनी जल्दिए प्रूशियन सरकार के सेंसरन के नजर में चढ़ गइल। जे अखबार में ऊ काम करत रहलन, ओकरा के ओकर बेबाक राजनीतिक टिप्पणी खातिर बंद करवा दिहल गइल, आ तंग आ के निराश मार्क्स पेरिस खातिर रवाना हो गइलन।

1843 के पेरिस ओह समय पूरे यूरोप के सामाजिक, राजनीतिक आ कलात्मक गतिविधि के सबसे बड़ केंद्र रहे, जहाँ यूरोप भर के क्रांतिकारी, समाजवादी आ विद्रोही लोग जुटल रहलन। पेरिस में मार्क्स प्रूडॉन आ बकुनीन जइसन नामी समाजवादियन आ विचारन संगे जुड़लन। बाकिर पेरिस में उनकर सबसे महत्वपूर्ण मिलाप फ्रेडरिक एंगेल्स संगे भइल। एंगेल्स इंग्लैंड के एगो अमीर कपड़ा कारखाना मालिक के बेटा रहलन, जे अपना बाप के कारखानन में मजदूरन के नारकीय आ बदतर हालत देख के खुद समाजवादी बन चुकल रहलन। मार्क्स आ एंगेल्स एक संगे मिल के ओह विचारन के गढ़े लगलन जे आगे चल के "क्रांतिकारी सर्वहारा समाजवाद" या जेकरा के आज दुनिया "कम्युनिज्म" (साम्यवाद) के नाम से जानेले, ओकर आधार बनल। आखिरकार, राजा-विरोधी आ उग्र लेखनी के कारण प्रूशियन सरकार के दबाव में 1845 में मार्क्स के फ्रांस से देस-निकाला दे दिहल गइल।

पेरिस छोड़े के बाद मार्क्स बेल्जियम पहुँचलन, जहाँ ऊ कारीगरन आ कामगारन के संगठन "कम्युनिस्ट लीग" से जुड़लन। 1847 में कम्युनिस्ट लीग मार्क्स आ एंगेल्स के आपन उद्देश्य, सिद्धांत आ कार्यक्रम के एगो घोषणापत्र लिखे के जिम्मेदारी दिहलस। ईहे रचना "कम्युनिस्ट घोषणापत्र" बनल, जेकरा मार्क्स 1848 में यूरोप भर में उठत क्रांति के लहर के देखत बड़े उत्साह आ तेज़ी संगे लिखलन। जब 1848 में जर्मनी में क्रांति के आगी सुलगल, त मार्क्स क्रांति के बढ़ावा देवे खातिर राइनलैंड गइलन। बाकिर जब क्रांति दबा दिहल गइल, त मार्क्स वापस पेरिस आ गइलन आ ओकरा कुछवे दिन बाद लंदन चले गइलन, जहाँ ऊ आपन बाकी पूरा जिनगी बितवलन।

लंदन में मार्क्स एह आस में बइठल रहलन कि क्रांति के आगी फेर से भभकी। एह तैयारी में ऊ यूरोप महाद्वीप के क्रांतिकारी नेता लोग संगे लगातार खत-पतरा आ संपर्क बनावले रहलन। ओह घड़ी ऊ इंग्लैंड के 'चार्टिस्ट' आ 'ट्रेड यूनियन' आंदोलनकारियन के नजरअंदाज कइलन, काहे से कि उनकरा लगत रहे कि ई लोग सादा सोच वाला आ निष्प्रभावी बा। समय बितला पर मार्क्स के अहसास भइल कि क्रांति हाल-फिलहाल तुरंत आवे वाला नइखे। एह से ऊ संगठनन के रोज़-रोज के झंझट से अलग हो के खुद के ब्रिटिश म्यूजियम (पुस्तकालय) के अध्ययन में गाड़ दिहलन, ताकि वर्ग-संघर्ष के ऐतिहासिक आधार पर गहिर शोध कर सकें। एह बरसन के कड़ा मेहनत आ शोध के फल मार्क्स के महान ग्रंथ 'दास कैपिटल' के रूप में सामने आइल, जेकर पहिला खंड 1867 में प्रकाशित भइल।

1863 में मार्क्स खातिर माहौल फेर से बदले लागल, जब फ्रांसीसी मजदूर लोग इंग्लैंड आइल ताकि ऊ लोग आर्थिक आ पूँजीवादी व्यवस्था के उखाड़ फेंके खातिर अंतरराष्ट्रीय मजदूर संगठन (फेडरेशन) बना सके। भले मार्क्स ओह समूह में शामिल कई विचारधारान से असहमत रहलन, बाकिर ऊ एह ऐतिहासिक घटना के महत्व के समझलन आ अपना खुद के एकांतवास के छोड़ के ओह आंदोलन में शामिल हो गइलन। मार्क्स चतुरता आ काबिलियत संगे एह संगठन के नेतृत्व में आपन जगह बनवलन, जेकरा के 'इंटरनेशनल' (प्रथम इंटरनेशनल) कहल गइल। ऊ प्रथम इंटरनेशनल के आपन ऐतिहासिक उद्घाटन भाषण दिहलन, जे उनकर सिद्धांतन के विजय-घोष रहे। आखिरकार मार्क्स के ऊ साधन मिल गइल जेकर आस ऊ 1847 से देखत रहलन—ऊ एगो अइसन अंतरराष्ट्रीय समाजवादी आंदोलन के बौद्धिक बुनियाद तैयार कर देले रहलन, जेकर वैचारिक कमान पूरा-पूरी उनकरे हाथ में रहे।

बाकिर मार्क्स के ई संतोष जियादा दिन ना टिक पावल। 1871 के 'पेरिस कम्यून', जहाँ मजदूर वर्ग इतिहास में पहली बेर खुद सत्ता पर काबिज भइल रहे, एक खूनी त्रासदी आ बर्बादी में खतम हो गइल। एह हिंसा से डर के इंग्लैंड के शांतिप्रिय मजदूर पीछे हट गइलन, आ फ्रांसीसी आंदोलन आपसी गुटबाज़ी में उलझ गइल। दूसरी ओर बकुनीन के अराजकतावादी (अनार्किस्ट) समर्थक लोग इंटरनेशनल के कमान मार्क्स के हाथ से छीने के पुरजोर कोशिश कइलस। मार्क्स आ अराजकतावादियन के बीच ई भीषण सैद्धांतिक लड़ाई अंततः 1876 में इंटरनेशनल के बिखरला पर खतम भइल।

आपन जिनगी के आखिरी कुछ बरसन में मार्क्स कवनो बड़ नया ग्रंथ ना लिख पइलन। बाकिर इंटरनेशनल के पूर्व सर्वमान्य नेता के रूप में उनकर साख अइसन रहे कि पूरा यूरोप, आ खास करके रूस के नया क्रांतिकारी समूह उनकरा से सलाह आ मार्गदर्शन माँगे खातिर आतुर रहत रहे। मार्क्स जहाँ तक संभव रहे एह नया विचारन आ नेता लोग के मदद कइलन, बाकिर फेर कबो कवनो आंदोलन के सीधा कमान ना सम्हाले आ गइलन। 1883 में लंदन में मार्क्स अंतिम साँस लिहलन, आ मरे दम तक ऊ ओही अवश्यंभावी क्रांति के राह देखत रहलन जेकर भविष्यवाणी ऊ अपना जीवन भर कइले रहलन।


\subsection*{भूमिका}
कम्युनिस्ट घोषणापत्र (मूल नाम: कम्युनिस्ट पार्टी के घोषणापत्र) राजनीतिक दार्शनिक कार्ल मार्क्स आ फ्रेडरिक एंगेल्स द्वारा 1848 में लिखल एगो छोट बाकिर अत्यंत प्रभावशाली रचना ह। तब से लेके आज तक ई दुनिया के सबसे महान आ असरदार राजनीतिक पांडुलिपियन में गिनल जाला। कम्युनिस्ट लीग के आदेश पर तैयार ई दस्तावेज, ओह लीग के मूल मकसद आ कार्यक्रम के दुनिया के सामने रखेला। ई साम्यवाद के कवनो काल्पनिक भविष्य के सपना देखावे के बजाय, इतिहास आ मउजूदा समय में चलत 'वर्ग-संघर्ष' आ पूँजीवाद के विकराल समस्यान के वैज्ञानिक आ विश्लेषणात्मक पड़ताल पेश करेला।

एह किताब में समाज आ राजनीति के स्वभाव पर मार्क्स आ एंगेल्स के कालजयी सिद्धांत बाड़ें, जिनकर कहनाम बा कि: "अब तक के सभन समाजन के इतिहास, वर्ग-संघर्षन के इतिहास रहल बा।" एकरा संगे-संगे, ई एह बात के भी रूपरेखा पेश करेला कि कइसे ओह समय के पूँजीवादी समाज अंततः समाजवाद में आ आगे चल के साम्यवाद (कम्युनिज्म) में बदल जाई।

\cleardoublepage
\section*{कम्युनिस्ट पार्टी के घोषणापत्र}
\addcontentsline{toc}{section}{कम्युनिस्ट पार्टी के घोषणापत्र}
[1888 के अंग्रेजी संस्करण से, फ्रेडरिक एंगेल्स द्वारा संपादित]

पूरा यूरोप पर एगो साया मँडरात बा—कम्युनिज्म (साम्यवाद) के साया। पुरनका यूरोप के सभन बड़का ताकत एह साया के भगावे आ ओकर तंत्र-मंत्र काटे खातिर एगो पबितर गठजोड़ (पवित्र गठबंधन) बना लिहले बाड़ी सँ: पोप आ जार, मेटरनिख आ गुइज़ो, फ्रांस के उग्रवादी आ जर्मनी के पुलिस-जासूस।

कवन अइसन विपक्षी पार्टी बा जेकरा के सत्ता में बइठल विरोधी लोग कम्युनिस्ट कह के कलंकित ना कइले होखे? कवन अइसन विपक्षी दल बा जे आपन से बेसी प्रगतिशील विपक्षी दलन पर या आपन घोर प्रतिक्रियावादी दुश्मनन पर कम्युनिस्ट होखे के दगाबाज़ी वाला दाग ना लगवले होखे?

एह सच्चाई से दू गो मुख्य बात निकल के समोर आवेला:

भाग {\engfont I}. कम्युनिज्म के अब पूरे यूरोप के सभन बड़का सत्ता-ताकत आपन आप में एगो महा-ताकत मान चुकल बाड़ी सँ।

भाग {\engfont II}. अब ई एकदम सही समय आ गइल बा कि कम्युनिस्ट लोग पूरा दुनिया के सामने खुलेआम आपन विचार, आपन मकसद, आपन नीति-नीयत आ आंदोलन के असली रूप पेश करे, आ कम्युनिज्म के साया वाली एह नानी-दादी के डरावनी कहानी के जवाब पार्टी के आपन आधिकारिक घोषणापत्र से देवे।

एही मकसद से अलग-अलग देसन के कम्युनिस्ट लोग लंदन में एक संगे जुटलन, आ नीचे दिहल घोषणापत्र के खाका तैयार कइलन, जेकरा के अंग्रेजी, फ्रेंच, जर्मन, इतालवी, फ्लेमिश आ डेनिश भाखा में छपवा के जारी कइल जाई।

\cleardoublepage
\section*{भाग {\engfont I}. बुर्जुआ आ सर्वहारा (पूँजीपति आ मज़दूर)}
\addcontentsline{toc}{section}{भाग {\engfont I}. बुर्जुआ आ सर्वहारा (पूँजीपति आ मज़दूर)}
अब तक के सभन समाजन के इतिहास, वर्ग-संघर्षन के इतिहास रहल बा।

स्वतंत्र नागरिक आ गुलाम, कुलीन (अभिजात) आ आम-जनता, सामंत-मालिक आ भूमिदास (खेतीहर मज़दूर), गिल्ड-मालिक आ दिहाड़ी कामगार—संक्षेप में कहल जाय त, शोषक आ शोषित हमेशा एक-दूसरा के सामने आड़ी आ के खड़ा रहलन। एह लोगन के बीच अनवरत लड़ाई चलत रहल—कबो छुप के त कबो सरेआम खुल के। ई लड़ाई हर बेर या त पूरा समाज के एगो क्रांतिकारी नया रूप में पुनर्गठन पर खतम भइल, या फेर लड़े वाला दुनो वर्गन के आपसी बर्बादी आ बिनाश पर।

इतिहास के पुरनका जमाना में हमनी लगभग हर जगह समाज के तरह-तरह के सामाजिक स्तरन आ पद-सोपानन (रैंक) में जटिल रूप से बँटल पावत बानी। प्राचीन रोम में कुलीन, घुड़सवार योद्धा, आम नागरिक आ गुलाम रहलन; मध्यकाल में सामंती सरदार, उनकर अधीनस्थ वसाल, गिल्ड-मालिक, कामगार, चेला-शगिर्द आ भूमिदास रहलन; आ एह वर्गन के भीतर भी छोटे-छोटे ऊँच-नीच के कई गो दर्जा बनल रहे।

सामंती समाज के खंडहरन आ ढहैला से उपजल आधुनिक बुर्जुआ समाज वर्ग-विरोधन के खतम ना कइलस। ई पुरनका के जगह बस नया वर्ग, शोषण के नया तरीका आ लड़ाई के नया रूप बना के खड़ा कर दिहलस। बाकिर हमनी के एह युग—बुर्जुआ युग के—एगो खास विशेषता बा: ई वर्ग-विरोधन के एकदम सादा आ स्पष्ट बना दिहले बा। पूरा समाज तेज़ी से दुगो विरोधी खेमा में, दुगो बड़का आमने-सामने खड़ा वर्गन में बँटत जात बा: बुर्जुआ (पूँजीपति) आ सर्वहारा (मज़दूर)।

मध्यकाल के भूमिदासन में से शुरुआती शहरन के आज़ाद नागरिक (बर्गर) उपजलन। एह शहरी नागरिकन में से हीं बुर्जुआ वर्ग के पहला बीज आ तत्व अंकुरित भइलन।

अमेरिका के खोज आ अफ्रीका के उत्तमाशा अंतरीप के चक्कर काटे वाला समुद्री रास्ता मिलला से उभरत बुर्जुआ वर्ग खातिर नया मैदान खुल गइल। ईस्ट-इंडीज़ आ चीन के विशाल बाजार, अमेरिका के उपनिवेशीकरण, उपनिवेशन संगे फलत-फूलत व्यापार, आ विनिमय के साधनन आ माल-असबाब में भारी बढ़त—एह सभ बातन व्यापार, जहाजरानी आ उद्योग-धंधन के अइसन अभूतपूर्व तेज़ी दिहलस जे पहिले कड़ियो ना देखल गइल रहे। एह तरहे डगमगात सामंती समाज के भीतर क्रांतिकारी तत्वन के बड़ी तेज़ी से विकास भइल।

उत्पादन के पुरनका सामंती ढाँचा, जहाँ औद्योगिक उत्पादन पर बंद गिल्डन के एकाधिकार रहत रहे, अब नया बाजारन के बढ़त माँग पूरा करे खातिर नाकाफी साबित होखे लागल। ओकर जगह मैन्युफैक्चरिंग (हस्तशिल्प कारखाना) प्रणाली ले लेहलस। गिल्ड-मालिकन के दरकिनार कर के मैन्युफैक्चरिंग मध्यम वर्ग आगे आइल; अलग-अलग कॉरपोरेट गिल्डन के बीच के श्रम-विभाजन हर एक कारखाने के भीतर के श्रम-विभाजन के सामने गायब हो गइल।

एह बीच बाजार लगातार बढ़त गइल आ माँग नित-नया चढ़ाव लेत रहल। अब मैन्युफैक्चरिंग प्रणाली भी माँग पूरा करे में असमरथ हो गइल। तब भाप आ मशीनरी उद्योग के दुनिया में महा-क्रांति ले आइल। मैन्युफैक्चरिंग कारखानन के जगह विशाल आधुनिक उद्योग (कारखाना प्रणाली) ले लेहलस; औद्योगिक मध्यम वर्ग के जगह औद्योगिक करोड़पतीयन, यानी पूरा-पूरा औद्योगिक सेना के कमान सम्हाले वाला आधुनिक बुर्जुआ (पूँजीपति) लोग ले लेहलस।

आधुनिक उद्योग दुनिया भर के वैश्विक बाजार (वर्ल्ड मार्केट) स्थापित कइलस, जेकर रास्ता अमेरिका के खोज तैयार कइले रहे। एह बाजार व्यापार, जहाजरानी आ जल-थल के यातायात साधने के अपार बढ़ावा दिहलस। एह विकास के असर फेर से उद्योग के विस्तार पर पड़ल; आ जइसे-जइसे उद्योग, व्यापार, जहाजरानी आ रेल-लाइन के पसार भइल, ओही अनुपात में बुर्जुआ वर्ग बढ़ल, ओकर पूँजी फलल-फूलल, आ मध्यकाल से चलल आवत हर पुरनका वर्ग के ऊ पीछे धकेल के पृष्ठभूमि में ढकेल दिहलस।

एह तरहे हमनी देखत बानी कि आधुनिक बुर्जुआ खुद एगो लंबे ऐतिहासिक विकास, उत्पादन आ विनिमय (लेन-देन) के तरीकन में भइल सिलसिलेवार क्रांतियन के उपज ह।

बुर्जुआ वर्ग के विकास के हर एक पड़ाव संगे ओह वर्ग के राजनीतिक प्रगति भी जुड़ल रहे। सामंती अभिजात वर्ग के हुकूमत में ई एगो सतावल आ दबावल वर्ग रहे; मध्यकालीन स्वायत्त नगरन (कम्यून) में ई एगो सशस्त्र आ स्वशासी संघ रहे; कहूँ स्वतंत्र शहरी गणराज्य (जइसे इटली आ जर्मनी में), त कहूँ राजशाही के कर-दाता "तीसरा दर्जा" (जइसे फ्रांस में); बाद में हस्तशिल्प (मैन्युफैक्चरिंग) काल में ई सामंती अभिजात वर्ग के बेअसर करे खातिर निरंकुश राजशाही के संतुलन-शक्ति आ असल में बड़-बड़ राजशाहियन के आधार-स्तंभ बनल। अंत में, आधुनिक उद्योग आ वैश्विक बाजार के स्थापना के बाद से, बुर्जुआ वर्ग आधुनिक प्रतिनिधि राज्य (डेमोक्रेटिक स्टेट) में पूरा-पूरी एकाधिकार वाला राजनीतिक हुकूमत जीत लेहलस। आधुनिक राज्य के कार्यपालिका (सरकार) आ कछहरी बस पूरा बुर्जुआ वर्ग के साझा मामला आ कारोबार के सम्हाले वाली एगो समिति (कमेटी) बन के रह गइल बा।

इतिहास के पन्ना में बुर्जुआ वर्ग बेहद क्रांतिकारी भूमिका निभावले बा।

जहाँ-जहाँ बुर्जुआ वर्ग सत्ता आ कमान सम्हालस, ओहिजा ऊ सभन सामंती, पितृसत्तात्मक आ भावनात्मक-भावुक संबंधन के छिन्न-भिन्न कर के अंत कर दिहलस। ऊ ओह रंग-बिरंग सामंती जंजीरन आ बन्धनन के बेदर्दी से फाड़ फेंकले बा जे आदमी के ओकर "प्राकृतिक आक़ा आ बड़हन लोग" से बान्ह के रखत रहे, आ आदमी आ आदमी के बीच नंगा स्वार्थ आ बेदर्द "नकदी के लेन-देन" (कैश पेमेंट) के सिवाय कवनो दोसर नाता ना छोड़े दिहलस। ऊ धार्मिक उन्माद, शूरवीरता के जोश आ तुच्छ भावुकता के सबसे पबितर आ आकाशीय आनंद के स्वार्थी आ बर्फीला पानी में डुबो के मार दिहलस। ऊ आदमी के व्यक्तिगत मोल आ मरजाद के विनिमय-मूल्य (दाम) में बदल दिहलस, आ अनगिनत प्रामाणिक आज़ादियन के जगह बस एक बेदर्दी आ अमानवीय आज़ादी स्थापित कइलस—मुफ्त व्यापार (फ्री ट्रेड)। एक शब्द में कहल जाय त, धर्म आ राजनीति के पर्दा में ढँकल शोषण के जगह ऊ नंगा, निर्लज्ज, सीधा आ क्रूर शोषण कायम कर दिहलस।

बुर्जुआ वर्ग एह से पहिले तक इज्जत आ श्रद्धा से देखल जाए वाला हर पेसा-रोजगार के ओकर पूजनीय मण्डली से बेदखल कर दिहलस। ऊ डॉक्टर, वकील, पादरी, कवि आ वैज्ञानिक सभकरा के आपन पगार पर काम करे वाला मजूर बना के रख दिहलस।

ई परिवार के चेहरा से ओकर भावुकता आ ममता के पर्दा नोच के फेंक दिहलस, आ पारिवारिक नाता के महज रुपिया-पैसा के नाता बना के छोड़ दिहलस।

ई साबित कर के देखा दिहलस कि मध्यकाल के जिस बर्बर आ बलशाली प्रदर्शन के प्रतिक्रियावादी लोग अतना तारीफ करेला, ओकर असली पूरक हिस्सा महज आलस आ कामचोरी रहे। ई सबसे पहिले देखावलस कि आदमी के मेहनत आ उद्यम का-का कारनामा कर सकता। ई मिस्र के पिरामिड, रोम के जल-प्रणाली (एक्वाडक्ट) आ गोथिक गिरजाघरन से भी कहीं बेसी महान आ अचरज भरे चमत्कार खड़ा कइलस; ई अइसन महा-अभियान चलावलस जेकर सामने इतिहास के पुरनका सभन देस-निकाला आ धर्मयुद्ध (क्रुसेड) फीका पड़ जालें।

बुर्जुआ वर्ग उत्पादन के साधनन (औजारन आ तकनीक) में लगातार क्रांति कइले बिना, आ एह तरह से उत्पादन के संबंधन आ ओकरा संगे-संगे पूरा सामाजिक संबंधन के बदले बिना जिंदा ना रह सकता। एकरा उलट, पुरनका उत्पादन तरीकन के जस के तस बचा के राखल पहिलका सभन औद्योगिक वर्गन के अस्तित्व के मुख्य शर्त रहे। उत्पादन में लगातार बदलाव आ क्रांति, सभ सामाजिक स्थितियन में अटूट उथल-पुथल, आ हमेशा बनल रहे वाला अनिश्चितता आ बेचैनी के माहौल बुर्जुआ युग के पहिलका सभन युगन से अलग करेला। सभ पुरनका जमावल आ बान्हल संबंध, आपन प्राचीन आ पूजनीय पूर्वाग्रह आ मान्यताओं संगे बह के गायब हो जालान; आ नया बनल संबंध सख्त होखे से पहिले ही पुरान आ बेकार हो जालान। जवन कुछ ठोस आ मजबूत बुझात रहे, ऊ सब भाप बन के हवा में उड़ गइल; जवन कुछ पबितर आ पूजनीय रहे, ओकर मान-मरजाद तार-तार हो गइल; आ अंत में आदमी मजबूर हो गइल कि ऊ आपन असली जीवन-स्थिति आ आपन संगी-साथियन संगे आपन संबंधन के खुली आ ठंडे दिमाग के आँख से देखे।

आपन बनल माल बिकावे खातिर लगातार बढ़त बाजार के जरूरत बुर्जुआ वर्ग के पूरे धरती पर खदेड़त फिरेले। ओकरा हर जगह बसे के पड़ेला, हर जगह जड़ जमावे के पड़ेला, आ हर जगह नया नाता-रिश्ता जोड़े के पड़ेला।

बुर्जुआ वर्ग वैश्विक बाजार के शोषण के जरिया हर देस के उत्पादन आ खपत के अंतरराष्ट्रीय (सर्वदेशीय) रूप दे दिहले बा। प्रतिक्रियावादियन के भारी कलेजा-काँप के बावजूद, ई उद्योग के गोड़ के नीचे से ओकर राष्ट्रीय जमीन खींच लेहले बा। सभ पुरनका राष्ट्रीय उद्योग नष्ट हो गइल बाड़ें या रोज-रोज नष्ट कइल जात बाड़ें। ओह लोग के जगह नया उद्योग लेत बाड़ें, जिनकर स्थापना सभन सभ्य देसन खातिर जिनगी-मउत के सवाल बन जाला—अइसन उद्योग जे अब देशी कच्चा माल पर ना, बलुक दूर-दराज के देसन से आइल कच्चा माल पर चलेलान; आ जिनकर बनल माल सिर्फ अपने देस में ना, बलुक दुनिया के हर कोना में खपत होला। पुरनका स्थानीय आ राष्ट्रीय अलगाव आ आत्मनिर्भरता के जगह अब हर दिशा में लेन-देन, देसन के एक-दूसरा पर सार्वभौमिक निर्भरता कायम हो गइल बा। जइसे भौतिक उत्पादन में, ओइसहीं बौद्धिक आ मानसिक उत्पादन में भी। अलग-अलग देसन के बौद्धिक रचना सभकर साझा संपत्ति बनत जात बा। राष्ट्रीय संकीर्णता आ संकुचित सोच अब असंभव होत जात बा, आ अनगिनत राष्ट्रीय आ स्थानीय साहित्यायन में से एगो "विश्व साहित्य" जन्म लेत बा।

बुर्जुआ वर्ग उत्पादन के सभ साधनन में तेज़ी से सुधार कर के, आ यातायात के साधनन के बेहद आसान बना के, सभन देसन के, इहँवा तक कि सबसे पिछड़ल देसन के भी सभ्यता के दायरा में खींच ले आवेला। ओकर सस्ता माल ओकर अइसन भारी तोप के गोला ह जेकरा से ऊ चीन के बड़का-बड़का दीवारन के ढा देवेला आ विदेशी लोग खातिर पिछड़ल देसन के जिद्दी नफरत के आत्मसमर्पण करे पर मजबूर कर देवेला। ई सभन देसन के, मिट जाए के डर दिखा के, बुर्जुआ उत्पादन प्रणाली अपनावे पर मजबूर करेला; ई ओह लोग के आपन कहल जाए वाली "सभ्यता" अपनावे पर मजबूर करेला, यानी खुद बुर्जुआ बन जाए खातिर मजबूर करेला। एक शब्द में कहल जाय त, ई पूरा दुनिया के आपन हीं रूप-रंग आ छवि में ढाल लेवेला।

ई गाँव आ देहात के शहरन के हुकूमत के अधीन कर दिहले बा। ई बड़-बड़ महानगरा बनवलस, गाँव के आबादी के मुकाबले शहरी आबादी के भारी बढ़ोतरी कइलस, आ एह तरहे आबादी के एगो बड़ हिस्सा के गाँव के बान्हल आ पिछड़ल जीवन के कूपमंडूकता से बाहर निकाल ले आइल। जइसे ई देहात के शहर पर निर्भर बनवलस, ओइसे ही ई पिछड़ल देसन के सभ्य देसन पर, किसान देसन के पूँजीपति देसन पर, आ पूरब (पूर्वी दुनिया) के पच्छिम (पश्चिमी दुनिया) पर निर्भर बना दिहलस।

बुर्जुआ वर्ग आबादी, उत्पादन के साधनन आ संपत्ति के बिखराव के लगातार खतम करत जात बा। ई उत्पादन के एक जगह जमा कर दिहले बा आ संपत्ति के कुछ गिने-चुने हाथन में केंद्रित कर दिहले बा। एकर लाज़मी नतीजा भइल राजनीतिक केंद्रीकरण। अलग-अलग हित, अलग कानून, अलग सरकार आ अलग टैक्स-प्रणाली वाला स्वतंत्र या ढीला जुड़ल इलाका सब समेट के एगो राष्ट्र में बाँध दिहल गइलन—एगो सरकार, एगो कानून, एगो राष्ट्रीय वर्ग-हित, एगो राष्ट्रीय सीमा आ एगो कस्टम-टैक्स सीमा। बुर्जुआ वर्ग आपन महज एक सौ साल के शासन में पहिलका सभन पीढ़ी से बेसी विशाल आ प्रचंड उत्पादक शक्तियन के खड़ा कर दिहलस। कुदरत के ताकतन पर आदमी के काबू, मशीनरी, उद्योग आ खेती में रसायन विज्ञान के इस्तेमाल, भाप-जहाज, रेल-लाइन, बिजली के तार, खेती खातिर पूरे महाद्वीपन के सफाई, नदियन के नहर बनावल, जमीन के नीचे से अचानक उपजल बड़-बड़ आबादी—पहिलका कवनो सदी खयालो ना कइले होखी कि सामाजिक श्रम के गोदी में अइसन महा-उत्पादक शक्तियाँ सूत रहल बाड़ी सँ!

एह से ई साफ देखात बा कि: जे उत्पादक साधनन आ विनिमय के आधार पर बुर्जुआ वर्ग खुद के खड़ा कइलस, ऊ सामंती समाज के भीतर ही तैयार भइल रहे। उत्पादन आ विनिमय के एह साधनन के विकास के एगो खास मोड़ पर, ऊ स्थितियाँ जिनकरा तहत सामंती समाज उत्पादन आ विनिमय करत रहे—खेती आ उद्योग के सामंती संगठन, यानी संपत्ति के सामंती संबंध—विकसित हो चुकल उत्पादक शक्तियन खातिर नाकाफी हो गइलन; ऊ सब भारी जंजीर आ रुकावट बन गइलन। ओह जंजीरन के टूटल तय रहे; आ ऊ तोड़ के फेंक दिहल गइलीं।

ओह जंजीरन के जगह आ गइल आज़ाद होड़ (मुक्त प्रतिस्पर्धा), आ ओकरा अनुकूल सामाजिक आ राजनीतिक व्यवस्था, संगे बुर्जुआ वर्ग के आर्थिक आ राजनीतिक हुकूमत।

आज ओही नियरा एगो आंदोलन हमनी के आँख के सामने चल रहल बा। आधुनिक बुर्जुआ समाज, जे उत्पादन, विनिमय आ संपत्ति के अइसन विशाल संबंधन के गढ़ले बा, ऊ ओह जादूगर नियर हो गइल बा जे तंत्र-मंत्र से पाताल के प्रेत-शक्तियन के बुला त लेहलस, बाकिर अब ओकरा पर काबू पावे में असमरथ बा। पिछला कई दहाई से उद्योग आ व्यापार के इतिहास बस आधुनिक उत्पादक शक्तियन के आधुनिक उत्पादन संबंधन के खिलाफ विद्रोह के इतिहास ह—ओह संपत्ति संबंधन के खिलाफ विद्रोह जे बुर्जुआ वर्ग आ ओकर शासन के वजूद के आधार बाड़ें। वाणिज्यिक संकटन (व्यापारिक मंदी) के नाम लेहल ही काफी बा, जे निश्चित अंतराल पर लौट-लौट के पूरे बुर्जुआ समाज के वजूद के गंभीर खतरा में डाल देवेलान। एह संकटन में ना खाली बनल-बनावल माल के बड़ हिस्सा, बलुक पहिले से तैयार उत्पादक शक्तियन के भी बड़ हिस्सा समय-समय पर तबाह कर दिहल जाला। एह संकटन में अइसन महामारी फाट जाले जे पहिलका सभन युगन में बेवकूफी आ अचरज लागित—अति-उत्पादन (ओवर-प्रोडक्शन) के महामारी। समाज अचानक से क्षणिक बर्बरता आ अकाल के हालत में ढकेल दिहल जाला; अइसन लागेला जइसे कवनो महा-अकाल या विनाशकारी युद्ध भोजन आ जीवन के सभ साधनन के काट के खतम कर दिहले होखे; उद्योग आ व्यापार तबाह देखात बाड़ें; आ काहे? काहे से कि समाज में ज़रूरत से बेसी सभ्यता, ज़रूरत से बेसी भोजन-सामग्री, ज़रूरत से बेसी उद्योग आ ज़रूरत से बेसी व्यापार हो गइल बा! समाज के पास मउजूदा उत्पादक शक्तियाँ अब बुर्जुआ संपत्ति के स्थितियन के आगे बढ़ावे में मदद ना करत बाड़ी सँ; ओकरा उलट, ऊ ओह स्थितियन खातिर ज़रूरत से बेसी ताक़तवर हो गइल बाड़ी सँ जिनकरा में ऊ बान्हल रहलें, आ जइसे ही ऊ एह जंजीरन के तोड़ेलीं, पूरे बुर्जुआ समाज में अव्यवस्था आ तबाही फैल जाले, आ बुर्जुआ संपत्ति के वजूद खतरा में पड़ जाला। बुर्जुआ समाज के दायरा अतना तंग बा कि ई आपन पैदा कइल अपार धन-संपदा के खुद में समावे में नाकाम हो जाला। आ बुर्जुआ वर्ग एह संकटन से कइसे उबरेला? एक ओर उत्पादक शक्तियन के बड़का हिस्सा के जबर्दस्ती तबाह कर के; दोसर ओर नया बाजारन पर कब्जा कर के आ पुरनका बाजारन के अउरी बेदर्दी से निचोड़ के। यानी कि, बेसी विशाल आ बेसी विनाशकारी संकटन के रास्ता साफ कर के, आ संकटन से बचे के साधनन के घटा के!

जे हथियाारन से बुर्जुआ वर्ग सामंतवाद के मार के गिरा देहले रहे, आज उहे हथियार खुद बुर्जुआ वर्ग के ओर तन गइल बाड़ें।

बाकिर बुर्जुआ वर्ग ना खाली आपन मउत के सामान आ हथियार गढ़ले बा; ई ओह हथियाारन के चलावे वाला लोगन के भी पैदा कर दिहले बा—आधुनिक कामगार वर्ग—सर्वहारा (मज़दूर)।

जइसे-जइसे बुर्जुआ वर्ग, यानी पूँजी, बढ़ेला, ओही अनुपात में सर्वहारा—आधुनिक मज़दूर वर्ग—के भी विस्तार होला। ई अइसन मज़दूरन के वर्ग ह जे तब तक ही जिंदा रह सकत बाड़ें जब तक ओह लोग के काम मिले, आ काम तब तक ही मिलेला जब तक ओह लोग के मेहनत पूँजी के बढ़ावत रहे। ई मज़दूर, जिनकरा खुद के टुकड़ा-टुकड़ा में बेचे के पड़ेला, व्यापार के कवनो दोसर सामान (माल) नियर बाड़ें, आ एह से ऊ बाज़ार के प्रतियोगिता के सभ उतार-चढ़ाव आ मंदी-तेज़ी के थपेड़न का शिकार बाड़ें।

मशीनन के भारी इस्तेमाल आ श्रम-विभाजन के कारण सर्वहारा के काम में कवनो व्यक्तिगत रंग या आकर्षण ना बचल। ऊ मशीन के महज एगो बेजान पुर्जा बन के रह जाला, आ ओकरा से बस सबसे सादा, एकरस आ आसानी से सीखल जाए वाला काम माँगल जाला। एह से मज़दूर के पैदा करे आ जिंदा राखे के खरचा लगभग ओकरा खातिर भोजन-कपड़ा आ ओकर नस्ल (बाल-बच्चन) के पाले तक सीमित हो के रह जाला। बाकिर कवनो सामान के दाम, आ एह से मेहनत (मजदूरी) के दाम भी ओकरा उत्पादन खरचा के बराबर होला। एही से, जइसे-जइसे काम के नीरसता आ तकलीफ बढ़ेला, ओइसे-ओइसे मजदूरी गिरत जाले। अउरी त अउरी, जइसे-जइसे मशीनन के इस्तेमाल आ श्रम-विभाजन बढ़ेला, ओही अनुपात में मेहनत के बोझ भी बढ़ेला—चाहे काम के घंटा बढ़ा के होखे, चाहे तय समय में बेसी काम ऐंठ के होखे, या मशीनन के रफ़्तार तेज़ कर के।

आधुनिक उद्योग पितृसत्तात्मक मालिक के छोट कारखाना के पूँजीपति के विशाल कारखाने में बदल दिहले बा। कारखाने में ठूँस के भरल मज़दूरन के झुंड के सैनिकन नियर संगठित कइल जाला। औद्योगिक सेना के साधारण सिपाही के रूप में ओह लोग के अधिकारी आ सार्जेंटन के सख्त कमान के नीचे राखल जाला। ऊ ना खाली बुर्जुआ वर्ग आ बुर्जुआ राज्य के गुलाम बाड़ें; बलुक ऊ रोज आ हर घड़ी मशीन के, ओवरसियर (सुपरवाइजर) के, आ सबसे बेसी खुद कारखाने के पूँजीपति मालिक के गुलाम बाड़ें। ई तानाशाही जितना बेबाकी से मुनाफ़ा कमाए के आपन एकमात्र मकसद घोषित करेले, ई अतना ही नीच, घिनौनी आ ज़हरीली लागत जाले।

हाथ के मेहनत में जितना कम हुनर आ ताक़त के ज़रूरत पड़ेला, यानी आधुनिक उद्योग जितना बेसी विकसित होला, ओतने ही पुरुषन के मेहनत के जगह मेहरारून (औरतन) के मेहनत ले लेवेले। उम्र आ लिंग के फरक मज़दूर वर्ग खातिर कवनो सामाजिक मायने ना रखे। सभन महज काम के औजार बाड़ें, जिनकर इस्तेमाल उनकर उम्र आ लिंग के हिसाब से कम या बेसी खरचीला होला।

कारखानेदार मालिक द्वारा मज़दूर के दबोच आ पेराई जइसहीं खतम होला आ ओकरा हाथ में नकद मजदूरी मिलेला, तइसहीं बुर्जुआ वर्ग के दोसर हिस्सा—मकान मालिक, दुकानदार, सूदखोर (महाजन)—ओकरा पर झपट पड़ेलान।

मध्यम वर्ग के नीचे के दर्जा—छोट व्यापारी, दुकानदार, अवकाशप्राप्त कारोबारी, कारीगर आ किसान—ई सब धीरे-धीरे खसक के सर्वहारा (मज़दूर वर्ग) में गिर जालें। एकर कारण ई बा कि उनकर छोट पूँजी आधुनिक उद्योग के बड़ पैमाने के मुकाबला ना कर पावे आ बड़े पूँजीपतियन के प्रतियोगिता में डूब जाले, आ दोसर कारण ई कि उनकर खास हुनर नया उत्पादन तरीकन के सामने बेकार हो जाला। एह तरहे सर्वहारा वर्ग में समाज के सभन वर्गन से लोग शामिल होत जालें।

सर्वहारा वर्ग विकास के कई गो दौरन से गुज़रेला। ओकर जनम संगे ही बुर्जुआ वर्ग संगे ओकर संघर्ष शुरू हो जाला। शुरू में ई लड़ाई एके-दुके मज़दूर लड़ेलान, फेर कारखाने के सभ कामगार, फेर कवनो एक इलाका में एकही धंधा के मज़दूर ऊ खास पूँजीपति मालिक के खिलाफ लड़ेलान जे सीधा ओह लोग के निचोड़त रहे। ऊ आपन हमला उत्पादन के बुर्जुआ संबंधन पर ना, बलुक खुद उत्पादन के औजारन (मशीनन) पर करेलान; ऊ मुकाबला करे वाला आयातित माल के तबाह कर देवेलान, मशीनन के तोड़-फोड़ कर देवेलान, कारखानन में आगी लगा देवेलान, आ मध्यकाल के कामगारन के खोइल दर्जा के जबर्दस्ती वापस पावे के कोशिश करेलान।

एह पड़ाव पर मज़दूर अबहियों पूरे देस में बिखरायल आ आपसी प्रतियोगिता से टूटल एगो असंगठित भीड़ बाड़ें। अगर ऊ कहूँ संगठित हो के जूटबो करेलान, त ई उनकर आपन सक्रिय एकता के नतीजा ना होला, बलुक बुर्जुआ वर्ग के साझिसी एकता के परिणाम होला। बुर्जुआ वर्ग आपन खुद के राजनीतिक मकसद पावे खातिर पूरे सर्वहारा के सड़क पर उतारे पर मजबूर होला आ कुछ समय खातिर अइसन करे में कामयाब भी हो जाला। एह से एह पड़ाव पर सर्वहारा आपन दुश्मनन संगे ना, बलुक आपन दुश्मनन के दुश्मनन संगे लड़ेला—निरंकुश राजशाही के अवशेषन संगे, जमींदारन संगे, गैर-औद्योगिक पूँजीपतियन आ छोट बुर्जुआ संगे। एह तरहे पूरा ऐतिहासिक आंदोलन बुर्जुआ वर्ग के हाथ में केंद्रित होला; आ एह तरहे हासिल हर जीत बुर्जुआ वर्ग के जीत बनेले।

बाकिर उद्योग के विकास संगे सर्वहारा वर्ग ना खाली गिनती में बढ़ेला; ई बेसी बड़का जमावड़ा में केंद्रित होला, ओकर ताक़त बढ़ेले आ ओकरा आपन एह ताक़त के अहसास होला। सर्वहारा के भीतर के अलग-अलग हित आ जीवन-स्थितियाँ बेसी से बेसी बराबर होत जालीं, काहे से कि मशीनरी काम के सभ अंतर के मिटा देवेले आ लगभग हर जगह मजदूरी के एकही नीचा स्तर पर ले आवेले। पूँजीपतियन के आपसी बढ़ती प्रतियोगिता आ ओकरा से आवे वाला वाणिज्यिक संकट मज़दूरन के मजदूरी के लगातार अनिश्चित बना देवेला। मशीनरी में लगातार सुधार आ तेज़ी से होत विकास ओह लोग के रोजी-रोटी के अउरी डाँवाडोल बना देवेला; एके-दुके मज़दूर आ एके-दुके पूँजीपति के बीच के टकराव अब दुगो वर्गन के बीच के टकराव के रूप लेवे लागेला। एह पर मज़दूर पूँजीपतियन के खिलाफ आपन संगठन (ट्रेड यूनियन) बनाए लगेलान; ऊ मजदूरी के दर बचावे खातिर एक संगे जुटेलान; ऊ एह समय-समय पर होखे वाला विद्रोह खातिर पहिले से तैयारी करे के खातिर स्थायी समिति बनवेलान। कहूँ-कहूँ ई लड़ाई उपद्रव आ बगावत के रूप ले लेवेले।

कबो-कबो मज़दूर जीतबो करेलान, बाकिर महज कुछ समय खातिर। उनकर लड़ायन के असली फल तत्काल मिले वाला नतीजा में ना, बलुक मज़दूरन के लगातार पसरत जाइब वाला एकता में छिपल बा। एह एकता के आधुनिक उद्योग द्वारा बनावल यातायात आ संचार के उन्नत साधन अउरी आसान बना देवेलान, जे अलग-अलग इलाका के मज़दूरन के एक-दूसरा से जोड़ देवेलान। एह संपर्क के ही ज़रूरत रहे ताकि एकही स्वभाव वाले अनगिनत स्थानीय संघर्षन के समेट के वर्ग-संघर्ष के रूप में राष्ट्रीय लड़ाई बना दिहल जाय। बाकिर हर वर्ग-संघर्ष एगो राजनीतिक संघर्ष होला। आ जिस एकता के हासिल करे में मध्यकाल के शहरी नागरिकन के आपन घटिया रास्तन के कारण सदियाँ लग गइल रहे, ओकरा के आधुनिक सर्वहारा रेल-लाइन आ संचार के बदौलत महज कुछ बरसन में हासिल कर लेवेलान।

सर्वहारा के एह तरहे वर्ग के रूप में, आ एह से राजनीतिक पार्टी के रूप में संगठन, मज़दूरन के आपसी प्रतियोगिता के कारण बार-बार टूटेला आ बिखर जाला। बाकिर ई हर बेर फेर से जी उठेला—अउरी मजबूत, अउरी अडिग, अउरी ताक़तवर हो के! ई बुर्जुआ वर्ग के आपसी फूट के फायदा उठा के मज़दूरन के खास हितन के कानूनी मान्यता दिलावे पर मजबूर करेला। एह तरहे इंग्लैंड में दस-घंटा के काम के कानून पास भइल रहे।

कुल मिला के पुरनका समाज के वर्गन के बीच के आपसी टकराव सर्वहारा के विकास के राह के कई तरह से आगे बढ़ावेला। बुर्जुआ वर्ग खुद के लगातार लड़ाई में घेरायल पावेला—पहिले अभिजात वर्ग (सामंतों) संगे; बाद में खुद बुर्जुआ वर्ग के ओह हिस्सन संगे जिनकर हित उद्योग के प्रगति के खिलाफ हो जालान; आ हमेशा विदेशी देसन के बुर्जुआ वर्ग संगे। एह सभन लड़ायन में ई खुद के सर्वहारा से गुहार लगावे, ओकर मदद माँगे आ एह तरहे ओकरा के राजनीति के मैदान में खींचे पर मजबूर पावेला। एह तरहे खुद बुर्जुआ वर्ग सर्वहारा के आपन राजनीतिक आ बौद्धिक शिक्षा के साधन मुहैया करावेला, यानी ऊ सर्वहारा के आपन हीं खिलाफ लड़े के हथियार सौंप देवेला।

अउरी त अउरी, जइसे हमनी देख चुकल बानी, उद्योग के प्रगति संगे सत्तारूढ़ वर्गन के पूरा-पूरा हिस्सा सर्वहारा में धकेल दिहल जालान, या कम से कम उनकर अस्तित्व के स्थितियन पर खतरा मँडरावे लागेला। ई लोग भी सर्वहारा वर्ग के ज्ञान आ प्रगति के नया तत्व प्रदान करेलान।

अंत में, ओह समय जब वर्ग-संघर्ष आपन निर्णायक घड़ी के नज़दीक पहुँचेला, त सत्तारूढ़ वर्ग के भीतर, आ असल में पूरा समाज के भीतर चलत विघटन के प्रक्रिया अतना उग्र आ साफ़ रूप ले लेवेले कि सत्तारूढ़ वर्ग के एगो छोट हिस्सा खुद के अलग कर के क्रांतिकारी वर्ग संगे जा मिलेला—ओह वर्ग संगे जेकरा हाथ में भविष्य बा। एह तरहे जइसे पहिलका दौर में अभिजात वर्ग के एगो हिस्सा बुर्जुआ वर्ग संगे जा मिलल रहे, ओइसहीं अब बुर्जुआ वर्ग के एगो हिस्सा सर्वहारा संगे जा मिलेला—खास करके बुर्जुआ विचारकों के ऊ हिस्सा जे पूरा ऐतिहासिक आंदोलन के सैद्धांतिक रूप से समझे के स्तर पर पहुँच चुकल बा।

आज बुर्जुआ वर्ग के सामने खड़ा सभन वर्गन में से केवल सर्वहारा ही असली मायने में क्रांतिकारी वर्ग ह। दोसर सभ वर्ग आधुनिक उद्योग के सामने सड़ के खतम हो जालान; सर्वहारा ओकर खास आ अनिवार्य उपज ह।

छोट मध्यम वर्ग, छोट कारखानेदार, दुकानदार, कारीगर, किसान—ई सब बुर्जुआ वर्ग से लड़ेलान ताकि मध्यम वर्ग के हिस्सा के रूप में आपन वजूद बचा सकें। एह से ऊ लोग क्रांतिकारी ना, बलुक रूढ़िवादी (संरक्षणवादी) बाड़ें। अउरी त अउरी, ऊ लोग प्रतिक्रियावादी बाड़ें, काहे से कि ऊ इतिहास के चक्का के पीछे घुमावे के कोशिश करेलान। अगर ऊ कबो क्रांतिकारी होबो करेलान, त महज एह डर से कि ऊ जल्दिए सर्वहारा में शामिल हो जाइहें; एह तरहे ऊ आपन मउजूदा ना, बलुक आपन भविष्य के हितन के रक्षा करेलान, ऊ आपन नजरिया छोड़ के सर्वहारा के नजरिया अपना लेवेलान।

"खतरनाक वर्ग" (लुंपेन-प्रोलिटेरिएट), सामाजिक कचरा, पुरनका समाज के सबसे निचला सड़ल परतन से फेकल गइल ई सड़त आ बेजान जमावड़ा, कबो-कबो सर्वहारा क्रांति द्वारा आंदोलन में बह के आ जाला; बाकिर ओकर जीवन-स्थितियाँ ओकरा के प्रतिक्रियावादी साज़िश के बिकाऊ औजार बने खातिर बेसी तैयार करेलीं।

सर्वहारा के स्थितियन में पुरनका समाज के स्थितियाँ पहले ही लगभग डूब के खतम हो चुकल बाड़ी सँ। सर्वहारा के पास कवनो संपत्ति नइखे; ओकर आपन मेहरारू आ लइकन संगे नाता बुर्जुआ पारिवारिक-नाता से कवनो मेल ना रखे; आधुनिक औद्योगिक श्रम, पूँजी के सामने आधुनिक दासता—चाहे इंग्लैंड में होखे या फ्रांस में, अमेरिका में होखे या जर्मनी में—ओकरा से राष्ट्रीय चरित्र के हर नामो-निशान पोछ दिहले बा। कानून, नैतिकता, धर्म ओकरा खातिर महज बुर्जुआ पूर्वाग्रह बाड़ें, जिनकरा पीछे उतने ही बुर्जुआ हित छुप के घात लगावले बइठल बाड़ें।

पहिलका सभन वर्ग जे सत्ता हासिल कइले रहलन, पूरे समाज के आपन हड़पे के स्थितियन के अधीन कर के आपन हासिल दर्जा के मजबूत करे के कोशिश कइलन। सर्वहारा समाज के उत्पादक शक्तियन के मालिक तब तक ना बन सकत, जब तक कि ऊ आपन पुरनका हड़पे के तरीका आ ओकरा संगे-संगे हर पुरनका हड़पे के तरीका के खतम ना कर दे। ओकरा पास आपन सुरक्षा आ मजबूती खातिर कवनो संपत्ति नइखे; ओकर एकमात्र मकसद व्यक्तिगत संपत्ति के सभन पुरनका सुरक्षा-कवचन आ बीमान के खतम कर देहल ह।

इतिहास के सभन पहिलका आंदोलन अल्पसंख्यकन के आंदोलन रहलन, या अल्पसंख्यकन के हित खातिर चलावल गइल रहलन। सर्वहारा आंदोलन विशाल बहुसंख्यक के हित में चलावल जाए वाला विशाल बहुसंख्यक के सचेत आ आज़ाद आंदोलन ह। सर्वहारा, हमनी के मउजूदा समाज के सबसे निचला दर्जा, तब तक हिल ना सकत, तब तक ऊपर ना उठ सकत, जब तक कि ओकरा ऊपर बइठल आधिकारिक समाज के सभन परतन के हवा में उड़ा ना दिहल जाय।

रूप में भले ना होखे, बाकिर तत्व में बुर्जुआ वर्ग संगे सर्वहारा के लड़ाई शुरू में एगो राष्ट्रीय लड़ाई होला। हर देस के सर्वहारा के सबसे पहिले आपन देस के बुर्जुआ वर्ग संगे हिसाब चुकता करे के पड़ेला।

सर्वहारा के विकास के सबसे आम चरणों के रेखांकित करत, हमनी ओह छुपल गृहयुद्ध के देखनी जे मउजूदा समाज के भीतर तब तक चलत रहे जब तक कि ऊ गृहयुद्ध खुल के क्रांति में ना बदल गइल, आ जहाँ बुर्जुआ वर्ग के हिंसक तख्तापलट सर्वहारा के शासन के बुनियाद रखलस।

अब तक के समाज के हर रूप, जइसे हमनी देख चुकल बानी, सतावे वाला आ सतावल जाए वाला वर्गन के विरोध पर आधारित रहल बा। बाकिर कवनो वर्ग के सतावे खातिर ई ज़रूरी बा कि ओकरा खातिर अइसन स्थितियाँ सुरक्षित कइल जाव जिनकरा में ऊ कम से कम आपन गुलामी वाली जिनगी बिता सके। सामंतवाद के दौर में भूमिदास खुद के कम्यून के सदस्य के रूप में ऊपर उठवलस, ओइसे ही निरंकुश सामंतवाद के जुआ नीचे छोट बुर्जुआ विकसित हो के बुर्जुआ बन गइल। एकरा उलट, आधुनिक मज़दूर उद्योग के प्रगति संगे ऊपर उठे के बजाय आपन वर्ग के अस्तित्व के स्थितियन से अउरी नीचे गिरत जाला। ऊ कँगाल (पॉपर) बन जाला, आ कंगाली आबादी आ धन से बेसी तेज़ी से फैलेला। एह से ई एकदम साफ़ हो जाला कि बुर्जुआ वर्ग अब समाज में सत्तारूढ़ वर्ग बने आ आपन अस्तित्व के स्थितियन के समाज पर सर्वोपरि कानून के रूप में थोपे खातिर अयोग्य हो चुकल बा। ई शासन करे खातिर एह से अयोग्य बा काहे से कि ई आपन गुलाम के ओकर गुलामी में भी जीवन-यापन के सुरक्षा देवे में असमरथ बा, काहे से कि ई ओकरा के अइसन हालत में गिरे से ना रोक पावत जेकरा में ओकरा के मज़दूर से पोषाए के जगह खुद मज़दूर के पोषे के पड़ेला। समाज अब एह बुर्जुआ वर्ग के हुकूमत में अउरी ना जी सकत, यानी एकर वजूद अब समाज संगे मेल ना खात।

बुर्जुआ वर्ग के वजूद आ हुकूमत के सबसे जरूरी शर्त ह पूँजी के निर्माण आ ओकर बढ़ोतरी; पूँजी के शर्त ह मजदूरी पर आधारित श्रम (वेज़-लेबर)। मजदूरी पर आधारित श्रम पूरा-पूरी मज़दूरन के आपसी प्रतियोगिता पर टिकल बा। उद्योग के प्रगति, जेकर अनिच्छुक बढ़ावा देवे वाला खुद बुर्जुआ वर्ग ह, प्रतियोगिता से उपजल मज़दूरन के अलगाव के जगह संगठन से उपजल उनकर क्रांतिकारी एकजुटता ले आवेले। एह तरहे आधुनिक उद्योग के विकास बुर्जुआ वर्ग के गोड़ के नीचे से ऊ बुनियाद ही काट फेंकेला जिस पर ऊ उत्पादन करेला आ माल हड़पेला। एह से बुर्जुआ वर्ग सबसे बेसी जेकर उत्पादन करेला, ऊ ह ओकर आपन कब्र खोदे वाला (कब्र खोदइया)। ओकर पतन आ सर्वहारा के जीत दुनो बराबर रूप से अवश्यंभावी (तय) बा।

\cleardoublepage
\section*{भाग {\engfont II}. सर्वहारा आ कम्युनिस्ट (मज़दूर आ साम्यवादी)}
\addcontentsline{toc}{section}{भाग {\engfont II}. सर्वहारा आ कम्युनिस्ट (मज़दूर आ साम्यवादी)}
पूरा सर्वहारा (मज़दूर वर्ग) संगे कम्युनिस्ट लोग के का नाता आ संबंध बा?

कम्युनिस्ट लोग मज़दूर वर्ग के दोसर पार्टियन के खिलाफ कवनो अलग आ थलग पार्टी ना बनावेलान।

सर्वहारा वर्ग के कुल हित से अलग उनकर कवनो आपन व्यक्तिगत हित ना होला।

ऊ सर्वहारा आंदोलन के आपन कवनो कट्टर या संकीर्ण सिद्धांतन के साँचा में ढाले के कोशिश ना करेलान।

कम्युनिस्ट लोग दोसर मज़दूर पार्टियन से सिर्फ एह दू गो बातन में अलग आ स्पष्ट बाड़ें: {\engfont (1)} अलग-अलग देसन के सर्वहारा वर्ग के राष्ट्रीय लड़ायन में, ऊ कवनो देस या राष्ट्रीयता से परे हट के पूरे सर्वहारा वर्ग के साझा हितन के सामने ले आभेलान आ ओकरा पर जोर देवेलान। {\engfont (2)} बुर्जुआ वर्ग के खिलाफ मज़दूर वर्ग के लड़ाई के जे-जे दौर आ पड़ाव से गुज़रे के पड़ेला, ऊ हर जगह आ हर समय पूरा आंदोलन के कुल हितन के प्रतिनिधित्व करेलान।

एह से, कम्युनिस्ट लोग एक ओर व्यावहारिक रूप से हर देस के मज़दूर वर्ग के पार्टियन के सबसे बेसी आगे बढ़े वाला, दृढ़ आ अडिग हिस्सा बाड़ें जे दोसर सभकरा के आगे ढकेलेला; दोसर ओर सैद्धांतिक रूप से, उनकरा पास पूरे सर्वहारा के मुकाबले ई बढ़त बा कि ऊ सर्वहारा आंदोलन के दिशा, ओकर स्थितियाँ आ ओकर अंतिम आम नतीजन् के साफ़-साफ़ समझेलान।

कम्युनिस्ट लोग के तत्काल मकसद उहे ह जे दोसर सभन सर्वहारा पार्टियन के ह: सर्वहारा के एगो संगठित वर्ग के रूप में तैयार कइल, बुर्जुआ वर्ग के वर्चस्व के उखाड़ फेंकल, आ सर्वहारा द्वारा राजनीतिक सत्ता पर कब्जा कइल।

कम्युनिस्ट लोग के सैद्धांतिक निष्कर्ष कवनो अइसन विचारन या सिद्धांतन पर आधारित नइखन जेकरा कवनो तथाकथित विश्व-सुधारक गढ़ले या खोजले होखे। ऊ महज हमनी के आँख के सामने चलत मउजूदा वर्ग-संघर्ष आ ऐतिहासिक आंदोलन से उपजल असली संबंधन के आम भाखा में बयान करेलान। मउजूदा संपत्ति संबंधन के खतमी कवनो कम्युनिज्म के अकेल खास विशेषता नइखे।

अतीत के सभन संपत्ति संबंध ऐतिहासिक स्थितियन के बदले संगे लगातार ऐतिहासिक बदलाव के शिकार होत आइल बाड़ें।

उदाहरण खातिर, फ्रांसीसी क्रांति सामंती संपत्ति के खतम कर के ओकरा जगह बुर्जुआ संपत्ति स्थापित कइलस।

कम्युनिज्म के खास विशेषता सामान्य संपत्ति के खतमी ना, बलुक बुर्जुआ संपत्ति के खतमी ह। बाकिर आधुनिक बुर्जुआ निजी संपत्ति माल पैदा करे आ हड़पे के अइसन ढाँचा के अंतिम आ सबसे पूरा रूप ह, जे वर्ग-विरोधन पर, यानी गिने-चुने लोगन द्वारा बहुसंख्यक लोगन के शोषण पर टिकल बा।

एह मायने में, कम्युनिस्ट लोग के सिद्धांत के एक वाक्य में समेटल जा सकत बा: निजी संपत्ति के खतमी।

हमनी कम्युनिस्ट लोग पर ई इल्जाम लगावल जाला कि हमनी आदमी के आपन मेहनत से कमाए आ हासिल करे के ओह अधिकार के खतम करे के चाहत बानी जेकरा के सभन व्यक्तिगत आज़ादी, उद्यम आ आत्मनिर्भरता के बुनियाद कहल जाला।

कड़ा मेहनत से कमायल, आपन बल पर हासिल, आपन कमाई के संपत्ति! का रउआ सभन के मतलब छोट कारीगर या छोट किसान के संपत्ति से बा—अइसन संपत्ति जे बुर्जुआ रूप से पहिले के ह? ओकरा के खतम करे के कवनो ज़रूरत नइखे; उद्योग के विकास ओकरा के बड़ पैमाने पर पहले ही नष्ट कर दिहले बा आ रोज-रोज नष्ट करत जात बा।

या रउआ सभन के मतलब आधुनिक बुर्जुआ निजी संपत्ति से बा?

बाकिर का मजदूरी पर आधारित श्रम (वेज़-लेबर) मज़दूर खातिर कवनो संपत्ति पैदा करेला? एकदम ना! ई पूँजी पैदा करेला—अइसन संपत्ति जे मजदूरी श्रम के निचोड़ेले, आ जे तब तक ना बढ़ सकत जब तक कि नया शोषण खातिर नयका मजदूरी श्रम के पैदा ना कर दे। संपत्ति आपन मउजूदा रूप में पूँजी आ मजदूरी श्रम के आपसी विरोध पर टिकल बा। आईं एह विरोध के दुनो पक्षन के जाँच कइल जाव।

पूँजीपति होखे के मतलब ना खाली कवनो व्यक्तिगत दर्जा ह, बलुक उत्पादन में एगो सामाजिक दर्जा ह। पूँजी एगो सामूहिक उपज ह, आ समाज के अनगिनत सदस्यन के साझा काम से, आ असल में समाज के सभन सदस्यन के साझा प्रयास से ही ओकरा के चलावल जा सकत बा।

एह से पूँजी कवनो व्यक्तिगत ताक़त ना, ई एगो सामाजिक ताक़त ह।

एह से, जब पूँजी के साझा संपत्ति में, पूरे समाज के सदस्यन के संपत्ति में बदल दिहल जाई, त ओकरा से व्यक्तिगत संपत्ति सामाजिक संपत्ति में ना बदले। केवल संपत्ति के सामाजिक चरित्र बदल जाला। ई आपन वर्ग-चरित्र खो देवेला।

अब आईं मजदूरी श्रम पर विचार कइल जाव।

मजदूरी श्रम के औसत दाम ह न्यूनतम मजदूरी (मिनिमम वेज), यानी जीवन-यापन के ओतने सामग्री जे मज़दूर के मज़दूर के रूप में जिंदा राखे खातिर एकदम जरूरी होखे। एह से मज़दूर आपन मेहनत से जे हासिल करेला, ऊ महज ओकर जिनगी के घिसटत राखे आ ओकरा के जिंदा राखे खातिर काफी होला। हमनी मेहनत के उपज के एह व्यक्तिगत हासिल के खतम करे के बिलकुल इरादा ना रखिला, जे इंसान के जीवन बचावे आ चलावे खातिर कइल जाला आ जेकरा से दोसर के मेहनत पर हुकूमत चलावे खातिर कवनो बेसी बचत ना बचे। हमनी बस एह हासिल के नीच आ घिनौना रूप के खतम करे के चाहिला, जेकरा तहत मज़दूर महज पूँजी बढ़ावे खातिर जीएला आ ओकरा के तब तक ही जीए दिहल जाला जब तक सत्तारूढ़ वर्ग के हित ओकर माँग करे।

बुर्जुआ समाज में जियत मेहनत महज जमा भइल मेहनत (पूँजी) के बढ़ावे के साधन ह। कम्युनिस्ट समाज में जमा भइल मेहनत मज़दूर के जीवन के व्यापक, समृद्ध आ खुशहाल बनावे के साधन ह।

एह से बुर्जुआ समाज में अतीत मउजूदा समय पर हुकूमत करेला; कम्युनिस्ट समाज में मउजूदा समय अतीत पर हुकूमत करेला। बुर्जुआ समाज में पूँजी आज़ाद बा आ ओकर आपन व्यक्तिगत वजूद बा, जबकि जियत इंसान निर्भर बा आ ओकर कवनो व्यक्तिगत पहचान नइखे।

आ एह हालत के खतमी के बुर्जुआ लोग व्यक्तिगत पहचान आ आज़ादी के खतमी कहेला! आ ई सही भी बा! बुर्जुआ व्यक्तिगत पहचान, बुर्जुआ स्वतंत्रता आ बुर्जुआ आज़ादी के खतमी ही हमनी के असल मकसद ह।

आज़ादी के मतलब मउजूदा बुर्जुआ उत्पादन स्थितियन के तहत होला—मुफ्त व्यापार, आज़ाद खरीद-बिक्री!

बाकिर अगर खरीद-बिक्री ही खतम हो जाई, त आज़ाद खरीद-बिक्री भी गायब हो जाई। आज़ाद खरीद-बिक्री के ई बात, आ स्वतंत्रता के बारे में बुर्जुआ वर्ग के दोसर सभन "बहादुरी भरे शब्द", मध्यकाल के बान्हल व्यापार आ रोक-टोक के मुकाबले में भले कवनो अर्थ राखत होखें, बाकिर जब कम्युनिस्टों द्वारा खरीद-बिक्री, बुर्जुआ उत्पादन स्थितियन आ खुद बुर्जुआ वर्ग के खतमी के बात आवेला, त एह बातन के कवनो मतलब ना रह जाला।

रउआ लोग एह बात से हक्का-बक्का बानी कि हमनी निजी संपत्ति खतम करे के चाहत बानी। बाकिर रउआ सभन के मउजूदा समाज में, आबादी के नौ-बता-दस (90\%) हिस्सा खातिर निजी संपत्ति पहिले ही खतम हो चुकल बा; ओकर कुछ लोगन खातिर वजूद महज एह से बा काहे से कि ऊ ओह 90\% लोगन के हाथ में नइखे! एह से रउआ हमनी पर ई इल्जाम लगावत बानी कि हमनी संपत्ति के अइसन रूप के खतम करे के चाहत बानी जेकर वजूद के जरूरी शर्त ही समाज के विशाल बहुसंख्यक पास कवनो संपत्ति ना होहल ह!

एक शब्द में, रउआ हमनी पर आपन (बुर्जुआ के) संपत्ति खतम करे के इल्जाम लगावत बानी। एकदम सही बात ह; हमनी के ईहे इरादा बा!

जे घड़ी से मेहनत के पूँजी, रुपिया-पैसा या लगान (रेंट) में ना बदलल जा सके, यानी अइसन सामाजिक ताक़त में ना बदलल जा सके जेकरा पर एकाधिकार कइल जा सके—यानी जे घड़ी से व्यक्तिगत संपत्ति के बुर्जुआ संपत्ति (पूँजी) में ना बदलल जा सके, ओही घड़ी से रउआ कहे लगिला कि व्यक्तिगत पहचान गायब हो गइल!

एह से रउआ के ई कबूल करे के पड़ी कि "व्यक्ति" से रउआ सभन के मतलब बुर्जुआ पूँजीपति मालिक के सिवाय दोसर कवनो इंसान नइखे! आ एह इंसान के रास्ता से हटावल आ ओकर वजूद असंभव बनावल एकदम ज़रूरी बा।

कम्युनिज्म कवनो इंसान से समाज के उपज के हड़पे के ताक़त ना छीने; ई बस ओकरा से एह हड़पे के जरिया दोसर के मेहनत के आपन गुलाम बनावे के ताक़त छीन लेवेला।

ई एतराज उठावल जाला कि निजी संपत्ति के खतमी पर सभ काम-काज बंद हो जाई आ चारों ओर आलस छाई जाई।

एह हिसाब से त बुर्जुआ समाज के कब के आलस में सड़ के तबाह हो जाए के चाहिए रहे; काहे से कि एह समाज में जे काम करेला ऊ कुछ हासिल ना करे, आ जे सब कुछ हासिल करेला ऊ काम ना करे! एह पूरा एतराज के मतलब बस ईहे पुनरुक्ति ह कि: जब पूँजी ना रही त मजदूरी श्रम भी ना रही।

भौतिक उपज पैदा करे आ हड़पे के कम्युनिस्ट तरीका के खिलाफ उठाए गइल सभन एतराज, ओही तरहे बौद्धिक उपज पैदा करे आ हड़पे के कम्युनिस्ट तरीका के खिलाफ भी उठावल जालें। जइसे बुर्जुआ खातिर वर्ग-संपत्ति के खतमी खुद उत्पादन के खतमी ह, ओइसे ही ओकरा खातिर वर्ग-संस्कृति के खतमी सभन संस्कृति के खतमी ह।

ऊ संस्कृति, जेकरा खतम होखे पर ऊ छाती पीटत बा, समाज के विशाल बहुसंख्यक खातिर महज मशीन नियर काम करे के ट्रेनिंग ह।

बाकिर रउआ हमनी से तब तक बहस मत करीं जब तक कि रउआ बुर्जुआ संपत्ति के खतमी के आपन बुर्जुआ आज़ादी, संस्कृति, कानून आदि के पैमाना से नापत बानी। रउआ सभन के विचार खुद रउआ बुर्जुआ उत्पादन आ संपत्ति के स्थितियन के उपज बाड़ें, जइसे रउआ न्यायशास्त्र महज रउआ वर्ग के इच्छा ह जेकरा के सभकरा खातिर कानून बना दिहल गइल बा—अइसन इच्छा जेकर स्वभाव आ दिशा रउआ वर्ग के अस्तित्व के आर्थिक स्थितियन से तय होला।

ई स्वार्थी भूल जे रउआ सभन के आपन मउजूदा उत्पादन आ संपत्ति के रूप से उपजल सामाजिक संबंधन के (जे इतिहास के विकास में पैदा हो के गायब हो जालें) कुदरत आ बुद्धि के सनातन (अमर) कानून माने पर मजबूर करेले—ई भूल रउआ पहिलका सभन सत्तारूढ़ वर्गन संगे साझा करत बानी। जवन बात रउआ प्राचीन संपत्ति या सामंती संपत्ति के मामले में साफ़ देखत बानी, ओकरा के आपन बुर्जुआ संपत्ति के मामले में माने से इनकार करत बानी!

परिवार के खतमी! इहँवा तक कि सबसे उग्रवादी लोग भी कम्युनिस्टों के एह कुख्यात प्रस्ताव पर भड़क उठेलान!

मउजूदा परिवार—बुर्जुआ परिवार—किस आधार पर टिकल बा? पूँजी पर, निजी मुनाफ़ा पर! आपन पूरा विकसित रूप में ई परिवार केवल बुर्जुआ वर्ग के बीच वजूद राखेला। बाकिर एकर पूरक रूप सर्वहारा वर्ग के बीच परिवार के व्यावहारिक अभाव आ सरेआम वेश्यावृत्ति में मिलेला।

बुर्जुआ परिवार आपन पूरक के खतम होखे संगे अपने-आप खतम हो जाई, आ दुनो पूँजी के खतम होखे संगे खतम हो जाइहें।

का रउआ हमनी पर ई इल्जाम लगावत बानी कि हमनी महतारी-बाप द्वारा लइकन के शोषण रोके के चाहत बानी? एह गुनाह के हमनी सहर्ष कबूल करत बानी!

बाकिर रउआ कहब कि हमनी घर के शिक्षा के सामाजिक शिक्षा से बदल के सबसे पबितर नाता के खतम कर देत बानी।

आ रउआ शिक्षा! का ऊ भी सामाजिक नइखे, जे समाज के स्थितियन से, स्कूलन आदि के जरिया समाज के सीधा या तिरछा दखल से तय होला? कम्युनिस्टों शिक्षा में समाज के दखल के ईजाद ना कइले; ऊ बस एह दखल के चरित्र बदले के चाहेलान आ शिक्षा के सत्तारूढ़ वर्ग के प्रभाव से बचावे के चाहेलान।

परिवार आ शिक्षा के बारे में, महतारी-बाप आ लइका के पबितर नाता के बारे में बुर्जुआ वर्ग के ई लच्छेदार बात अउरी घिनौनी लागत जाले जब आधुनिक उद्योग द्वारा सर्वहारा के सभन पारिवारिक बन्धन फाड़ के फेंक दिहल जालें आ उनकर लइकन के महज व्यापारिक सामान आ काम के औजार बना दिहल जाला।

बाकिर रउआ कम्युनिस्ट लोग औरतन के साझी (कॉमन) बनावे के चाहत बानी—पूरा बुर्जुआ वर्ग एक सुर में चिल्ला उठेला!

बुर्जुआ आपन मेहरारू में महज उत्पादन के औजार देखेला। ऊ सुनेला कि उत्पादन के औजारन के साझा इस्तेमाल होई, त स्वाभाविक रूप से ऊ एह नतीजे पर पहुँचेला कि औरतन के भी सभकर साझा होखे के अइसन मउका आई!

ओकरा एह बात के अंदेशा तक नइखे कि असली मकसद औरतन के महज उत्पादन के औजार के दर्जा से आज़ाद करावल ह!

बाकी, औरतन के साझा बनावे के बात पर हमनी के बुर्जुआ वर्ग के ई पबितर आ धर्मनिष्ठ गुस्सा से बेसी हास्यास्पद दोसर कुछ ना हो सकत, जेकरा बारे में ऊ दावा करेलान कि ई कम्युनिस्टों द्वारा खुलेआम कानूनन लागू कइल जाई। कम्युनिस्टों के औरतन के साझी बनावे के कवनो ज़रूरत नइखे; ई अनादि काल से लगभग मउजूदा रहल बा।

हमनी के बुर्जुआ, सर्वहारा के मेहरारून आ बेटियन के आपन मर्जी पर रखे (संगे सरेआम वेश्यावृत्ति के बात छोड़ दीं) से संतोस ना पा के, एक-दूसरा के मेहरारून के बहकावे आ फँसावे में सबसे बेसी आनंद लेवेलान!

बुर्जुआ बिआह (शादी) असल में साझा औरतन के एगो ढाँचा ह, एह से कम्युनिस्टों पर बेसी से बेसी ई इल्जाम लगावल जा सकत बा कि ऊ छुपल आ ढोंगी साझी औरतन के जगह खुलेआम कानूनी साझी औरतन के लावे के चाहत बाड़ें। बाकी ई अपने-आप साफ़ बा कि मउजूदा उत्पादन प्रणाली के खतमी संगे ओकरा से उपजल औरतन के साझी रूप भी खतम हो जाई—यानी सरकारी आ निजी दुनो तरह के वेश्यावृत्ति खतम हो जाई।

कम्युनिस्टों पर आगे ई इल्जाम लगावल जाला कि ऊ देस आ राष्ट्रीयता खतम करे के चाहत बाड़ें।

मज़दूरन के कवनो देस ना होला। हमनी ओह लोग से ऊ ना छीन सकत जे उनकरा पास नइखे! काहे से कि सर्वहारा के सबसे पहिले राजनीतिक वर्चस्व हासिल करे के बा, देस के अग्रणी वर्ग के रूप में ऊपर उठे के बा, खुद के राष्ट्र के रूप में गठित करे के बा, एह से ऊ अब तक खुद राष्ट्रीय बा, भले बुर्जुआ मायने में ना होखे।

आधुनिक बुर्जुआ वर्ग के विकास, व्यापार के आज़ादी, वैश्विक बाज़ार, आ उत्पादन व जीवन-स्थितियन के एकरूपता के कारण देसन के बीच के राष्ट्रीय अंतर आ विरोध रोज-रोज गायब होत जात बाड़ें।

सर्वहारा के हुकूमत एह अंतरन के अउरी तेज़ी से गायब कर दी। कम से कम मुख्य सभ्य देसन के साझा कार्रवाई सर्वहारा के आज़ादी के पहिलका स्थितियन में से एक ह।

जइसे-जइसे एक इंसान द्वारा दोसर इंसान के शोषण खतम कइल जाई, ओइसे-ओइसे एक देस द्वारा दोसर देस के शोषण भी खतम कइल जाई। देस के भीतर वर्ग-विरोध खतम होखे संगे-संगे देसन के एक-दूसरा के प्रति दुश्मनी भी खतम हो जाई।

धार्मिक, दार्शनिक आ वैचारिक नजरिया से कम्युनिज्म पर लगावल गइल इल्जाम गंभीर जाँच के लायक नइखन।

का एह बात के समझे खातिर गहिर ज्ञान के ज़रूरत बा कि इंसान के विचार, नजरिया आ अवधारणा—एक शब्द में इंसान के चेतना—ओकर भौतिक वजूद के स्थितियन, ओकर सामाजिक संबंधन आ ओकर सामाजिक जीवन में भइल हर बदलाव संगे बदल जाले?

विचारन के इतिहास ई साबित करे के सिवाय दोसर का करेला कि बौद्धिक उत्पादन आपन चरित्र ओही अनुपात में बदलेला जिस अनुपात में भौतिक उत्पादन बदलेला? हर युग के सत्तारूढ़ विचार हमेशा ओकर सत्तारूढ़ वर्ग के विचार रहल बाड़ें।

जब लोग अइसन विचारन के बात करेलान जे समाज में क्रांति ले आवेलान, त ऊ बस एह सच्चाई के बयान करेलान कि पुरनका समाज के भीतर नया समाज के तत्व तैयार हो चुकल बाड़ें, आ पुरनका विचारन के बिखरना पुरनका जीवन-स्थितियन के बिखरे संगे ताल मिला के चलेला।

जब प्राचीन दुनिया आपन अंतिम साँस लेत रहे, त प्राचीन धर्मन के ईसाई धर्म हरा दिहले रहे। जब 18वीं सदी में ईसाई विचार विवेकवादी (तर्कवादी) विचारन के सामने घुटना टेक दिहले, त सामंती समाज ओह समय के क्रांतिकारी बुर्जुआ वर्ग संगे आपन अंतिम मउत के लड़ाई लड़त रहे। धार्मिक स्वतंत्रता आ ज़मीर के आज़ादी के विचार महज ज्ञान के क्षेत्र में मुफ़्त प्रतियोगिता के हुकूमत के बयान करत रहलन।

लोग कही: "बेशक, ऐतिहासिक विकास के दौरान धार्मिक, नैतिक, दार्शनिक आ कानूनी विचार बदले बाड़ें। बाकिर धर्म, नैतिकता, दर्शन, राजनीति आ कानून एह बदलावन के बावजूद हमेशा बचे रहल बाड़ें।"

"एकरा अलावा कुछ अमर (शाश्वत) सच बाड़ें, जइसे आज़ादी, न्याय आदि, जे समाज के सभन स्थितियन में साझा बाड़ें। बाकिर कम्युनिज्म अमर सचन के खतम करेला, ई धर्म आ नैतिकता के नया आधार पर खड़ा करे के बजाय खतम कर देवेला; एह से ई अतीत के सभन ऐतिहासिक अनुभवन के उलट काम करेला।"

ई इल्जाम किस बात पर आ के टिकेला? अतीत के सभन समाज के इतिहास वर्ग-विरोधन के विकास में बीतल बा—अइसन विरोध जे अलग-अलग युगन में अलग-अलग रूप लेत रहलन।

बाकिर चाहे ऊ कवनो रूप लेहले होखें, एक बात सभन पहिलका युगन में साझा बा—समाज के एक हिस्सा द्वारा दोसर हिस्सा के शोषण! एह से एह में कवनो अचरज नइखे कि पहिलका युगन के सामाजिक चेतना, आपन सभन विविधता के बावजूद, कुछ साझा रूपन या आम विचारन के भीतर घूमती रहे जे वर्ग-विरोधन के पूरा तरह गायब भइला बिना खतम ना हो सकत।

कम्युनिस्ट क्रांति पारंपरिक संपत्ति संबंधन संगे सबसे बेबाक आ गहिर नाता-तोड़ाव ह; एह से एह में कवनो अचरज नइखे कि एकर विकास पारंपरिक विचारन संगे सबसे गहिर नाता-तोड़ाव ले आवेला।

बाकिर आईं अब कम्युनिज्म पर बुर्जुआ एतराज़न के बात बंद कइल जाव।

हमनी ऊपर देख चुकल बानी कि मज़दूर वर्ग द्वारा क्रांति के पहिलका कदम सर्वहारा के सत्तारूढ़ वर्ग के दर्जा पर उठावल आ लोकतंत्र के लड़ाई जीतल ह।

सर्वहारा आपन राजनीतिक वर्चस्व के इस्तेमाल बुर्जुआ वर्ग से धीरे-धीरे पूरा पूँजी छीन लेवे खातिर करी, ताकि राज्य के हाथ में—यानी सत्तारूढ़ वर्ग के रूप में संगठित सर्वहारा के हाथ में—उत्पादन के सभन औजारन के केंद्रित कइल जा सके, आ कुल उत्पादक शक्तियन के जितना तेज़ी से हो सके बढ़ावल जा सके।

बेशक शुरू में ई काम संपत्ति के अधिकारन आ बुर्जुआ उत्पादन स्थितियन पर जबर्दस्ती आ निरंकुश हमला कइले बिना ना हो सकत; यानी अइसन उपायन के जरिया जे आर्थिक रूप से नाकाफी आ टिके योग्य ना लगें, बाकिर आंदोलन के दौरान ऊ खुद से आगे निकल जालें, पुरनका सामाजिक ढाँचा पर अउरी हमला ज़रूरी बना देवेलान, आ उत्पादन प्रणाली के पूरा तरह क्रांतिकारी बनावे के साधन के रूप में अनिवार्य बाड़ें।

ई उपाय अलग-अलग देसन में अलग-अलग होइहें।

फिर भी सबसे विकसित देसन में नीचे दिहल उपाय बड़ पैमाने पर लागू होइहें:


\subsection*{{\engfont 1}. जमीन पर निजी संपत्ति के खतमी आ लगान (रेंट) के जनहित आ सार्वजनिक कामन में इस्तेमाल।}

\subsection*{{\engfont 2}. भारी प्रगतिशील {\engfont (Graduated)} आयकर (इंकम टैक्स)।}

\subsection*{{\engfont 3}. उत्तराधिकार (बंश-संपत्ति हक) के पूरा-पूरी खतमी।}

\subsection*{4. देस छोड़ के भागे वालन आ विद्रोहियन के संपत्ति के ज़ब्ती।}

\subsection*{5. सरकारी बैंक आ सरकारी पूँजी के जरिया कर्ज़ आ साख (क्रेडिट) के राज्य के हाथ में एकाधिकार संगे केंद्रीकरण।}

\subsection*{6. यातायात आ संचार के सभन साधनन के राज्य के हाथ में केंद्रीकरण।}

\subsection*{7. सरकारी कारखानन आ उत्पादक औजारन के विस्तार; साझा योजना के तहत परती जमीन के खेती योग्य बनावल आ माटी के आम सुधार।}

\subsection*{8. सभन लोग खातिर काम करे के बराबर जिम्मेदारी; खास करके खेती खातिर औद्योगिक सेना के स्थापना।}

\subsection*{9. खेती आ उद्योग-धंधा के एक संगे मिलावल; पूरे देस में आबादी के बराबर बंटवारा कर के गाँव आ शहर के अंतर के धीरे-धीरे खतम कइल।}

\subsection*{10. सभन लइकन खातिर सरकारी स्कूलन में मुफ़्त शिक्षा; कारखानन में लइकन के काम करे के मउजूदा रूप के खतमी; शिक्षा के औद्योगिक उत्पादन संगे जोड़ल, आदि।}
जब विकास के दौरान वर्ग-भेद गायब हो जाइहें आ पूरा उत्पादन पूरा राष्ट्र के विशाल संघ के हाथ में केंद्रित हो जाई, त सार्वजनिक सत्ता आपन राजनीतिक चरित्र खो दी। असली मायने में राजनीतिक सत्ता महज एक वर्ग द्वारा दोसर वर्ग के सतावे खातिर संगठित ताक़त ह। अगर बुर्जुआ वर्ग संगे लड़ाई के दौरान सर्वहारा मजबूरी में खुद के वर्ग के रूप में संगठित करेला, अगर क्रांति द्वारा खुद के सत्तारूढ़ वर्ग बनावेला आ सत्तारूढ़ वर्ग के रूप में पुरनका उत्पादन स्थितियन के बलपूर्वक खतम कर देवेला, त ऊ एह स्थितियन संगे वर्ग-विरोधन आ आम वर्गन के अस्तित्व के स्थितियन के भी खतम कर दी, आ एह तरहे एक वर्ग के रूप में आपन हुकूमत भी खतम कर दी।

पुरनका बुर्जुआ समाज के जगह, ओकर वर्गन आ वर्ग-विरोधन के जगह, हमनी के पास अइसन संघ (समाज) होई जहाँ हर एक आदमी के आज़ाद विकास, सभन लोग के आज़ाद विकास के शर्त होई।

\cleardoublepage
\section*{भाग {\engfont III}. समाजवादी आ कम्युनिस्ट साहित्य}
\addcontentsline{toc}{section}{भाग {\engfont III}. समाजवादी आ कम्युनिस्ट साहित्य}

\subsection*{{\engfont 1}. प्रतिक्रियावादी समाजवाद}
\addcontentsline{toc}{subsection}{{\engfont 1}. प्रतिक्रियावादी समाजवाद}

\subsection*{{\engfont A}. सामंती समाजवाद}
आपन ऐतिहासिक स्थिति के कारण फ्रांस आ इंग्लैंड के सामंती अभिजात वर्ग के ई फर्ज बन गइल कि ऊ लोग आधुनिक बुर्जुआ समाज के खिलाफ तीखा पर्चा आ साहित्य लिखे। जुलाई 1830 के फ्रांसीसी क्रांति में आ इंग्लैंड के सुधार आंदोलन में ई सामंती अभिजात लोग फेर से ओह घमंडी आ नीच नव-धनाढ्य बुर्जुआ वर्ग के सामने हार गइलन। ओकरा बाद से कवनो गंभीर राजनीतिक लड़ाई के गुंजाइश ना रह गइल। केवल एगो साहित्यिक लड़ाई ही संभव बचल। बाकिर साहित्य के मैदान में भी पुरनका सामंती राज-वापसी काल के जय-जयकार वाली आवाज़ अब असंभव हो चुकल रहे।

सहानुभूति जगावे खातिर अभिजात वर्ग के मजबूरी हो गइल कि ऊ दिखावा खातिर आपन स्वार्थ भुला के बुर्जुआ वर्ग के खिलाफ आपन इल्जाम केवल शोषित मज़दूर वर्ग के हित में रखे। एह तरहे ऊ लोग नया मालिक पर व्यंग्य-गीत गा के आ ओकर कान में आवे वाला विनाश के ज़हरीली भविष्यवाणी फुसफुसा के आपन बदला लेवे लागल।

एह तरहे सामंती समाजवाद पैदा भइल: आधा शोक-गीत, आधा उपहास; आधा अतीत के गूँज, आधा भविष्य के धमकी; कबो-कबो ओकर तीखा, चतुर आ छुरा नियर चुभे वाला आलोचना बुर्जुआ के कलेजा के छेद देत, बाकिर आधुनिक इतिहास के रफ़्तार समझे में पूरा असमरथता के कारण ओकर असर हमेशा ही हास्यास्पद रहत।

अभिजात वर्ग, लोग के आपन ओर खींचे खातिर, मज़दूरन के भीख मँगइया झोली के झंडा बना के लहरावे। बाकिर जब-जब जनता उनकर साथ देवे खातिर भागे, तब-तब जनता उनकर पिछवाड़ा पर पुरनका सामंती कुल-चिन्ह देख लेवे, आ जोर-जोर से ठहाका मार के हँसत उनकरा से दूरी बना लेवे!

फ्रांस के 'लेजिटिमिस्ट' आ "यंग इंग्लैंड" के एगो हिस्सा ई तमाशा पेश कइलस।

ई बतवला में कि उनकर शोषण के तरीका बुर्जुआ से अलग रहे, सामंतवादी लोग ई भुला जात रहलन कि ऊ लोग जिन स्थितियन आ शर्तन के तहत शोषण करत रहे, ऊ बिलकुल अलग आ अब पुरान-बेकार हो चुकल रहलीं। ई देखावला में कि उनकर शासन के तहत आधुनिक मज़दूर वर्ग कबो मउजूदा ना रहे, ऊ ई भुला जात रहलन कि आधुनिक बुर्जुआ वर्ग खुद उनकरे समाज के लाज़मी औलाद ह!

बाकी, ऊ लोग आपन आलोचना के प्रतिक्रियावादी स्वभाव के छुपावे में अतना असमरथ रहलन कि बुर्जुआ वर्ग पर उनकर मुख्य इल्जाम बस ईहे रह गइल कि बुर्जुआ शासन के तहत अइसन वर्ग पैदा होत बा जे पुरनका सामाजिक ढाँचा के जड़-पेड़ से उखाड़ के फेंक दी।

जे बात पर ऊ बुर्जुआ के डाँटत रहलन, ऊ ई ना कि बुर्जुआ मज़दूर वर्ग बनावत बा, बलुक ई कि ऊ एक क्रांतिकारी सर्वहारा (मज़दूर वर्ग) बनावत बा!

एह से राजनीतिक व्यवहार में ऊ लोग मज़दूर वर्ग के खिलाफ सभन दमनकारी कदम में शामिल हो गइलन; आ आम ज़िंदगी में, आपन बड़ी-बड़ी बातन के बावजूद, ऊ उद्योग के पेड़ से गिरल सोना के सेब उठावे खातिर झुक जालें, आ सच, प्यार व इज़्ज़त सब बेच के ऊन, चुकंदर के चीनी आ आलू के शराब के धंधा करे लगेलान।

जइसे पादरी हमेशा जमींदार संगे कदम से कदम मिला के चलल बा, ओइसे ही पादरीवादी समाजवाद हमेशा सामंती समाजवाद संगे चलल बा।

ईसाई संयम आ वैराग्य के समाजवादी रंग देहल सबसे आसान काम ह। का ईसाई धर्म निजी संपत्ति, शादी आ राज्य के खिलाफ ना बोललस? का ई एह सभन के जगह दान-पुण्य आ गरीबी, ब्रह्मचर्य आ देह के तपस्या, मठ-जीवन आ माता चर्च के उपदेश ना दिहलस? ईसाई समाजवाद महज ऊ पबितर पानी ह जेकरा से पादरी अभिजात वर्ग के कलेजा के जलन आ ईर्ष्या पर छिड़काव करेला।


\subsection*{{\engfont B}. छोट-बुर्जुआ समाजवाद}
सामंती अभिजात वर्ग ही अकेला अइसन वर्ग ना रहे जेकरा के बुर्जुआ वर्ग तबाह कइले होखे, आ ना ही ऊ अकेला वर्ग रहे जेकर वजूद के स्थितियां आधुनिक बुर्जुआ समाज के हवा में सड़ के खतम हो गइलीं। मध्यकालीन नगरवासी आ छोट किसान-मालिक आधुनिक बुर्जुआ वर्ग के पूर्वज रहलन। जे देसन में औद्योगिक आ व्यापारिक विकास अबहीं तक कम भइल बा, ओहिजा ई दुनो वर्ग अबहीं तक उभरत बुर्जुआ संगे घिसटत जिंदा बाड़ें।

जे देसन में आधुनिक सभ्यता पूरा रूप से विकसित हो चुकल बा, ओहिजा एगो नया छोट-बुर्जुआ वर्ग बन गइल बा, जे सर्वहारा आ बुर्जुआ के बीच झूलत रहेला आ बुर्जुआ समाज के पूरक हिस्सा के रूप में खुद के नया रूप देत रहेला। बाकिर एह वर्ग के सदस्य प्रतियोगिता के दबाव में लगातार सर्वहारा में धकेलात जालें, आ जइसे-जइसे आधुनिक उद्योग बढ़ेला, ऊ लोग ओह घड़ी के नज़दीक आवत देखत बाड़ें जब ऊ आधुनिक समाज में एक आज़ाद हिस्सा के रूप में पूरा तरह गायब हो जाइहें, आ उद्योग, खेती व व्यापार में उनकर जगह सुपरवाइजर, कारिंदा आ दुकानदार ले लीहें।

फ्रांस जइसन देसन में, जहाँ किसान आबादी के आधा से काफी बेसी बाड़ें, ई स्वाभाविक रहे कि जे लेखक बुर्जुआ वर्ग के खिलाफ सर्वहारा के पक्ष लेत रहलन, ऊ बुर्जुआ शासन के आलोचना में किसान आ छोट-बुर्जुआ के पैमाना अपनावें, आ एह मध्यवर्ती वर्गन के नजरिया से मज़दूर वर्ग खातिर लड़ाई उठावें। एह तरहे छोट-बुर्जुआ समाजवाद पैदा भइल। सिसमॉन्डी एह स्कूल के सिरमौर रहलन, ना खाली फ्रांस में, बलुक इंग्लैंड में भी।

समाजवाद के एह स्कूल आधुनिक उत्पादन स्थितियन में छिपल विरोधाभासन के बड़ी तेज आ सटीक नजर से खोल के रख दिहलस। ऊ अर्थशास्त्रियन के ढोंगी बहाना उजागर कइलस। ऊ मशीनन आ श्रम-विभाजन के विनाशकारी असर, पूँजी आ जमीन के कुछ गिने-चुने हाथन में जमाव, अति-उत्पादन आ व्यापारिक संकट, छोट-बुर्जुआ आ किसानन के अनिवार्य बर्बादी, सर्वहारा के बदहाली, उत्पादन में अव्यवस्था, धन के बंटवारा में चीखत असमानता, देसन के बीच विनाशकारी औद्योगिक युद्ध, पुरनका नैतिक बंधन, पुरनका पारिवारिक संबंध आ पुरनका राष्ट्रीयता के बिखराव—ई सब साफ़-साफ़ देखा दिहलस।

बाकिर आपन सकारात्मक मकसद में ई समाजवाद या त पुरनका उत्पादन आ विनिमय के साधन, ओकरा संगे पुरनका संपत्ति संबंध आ पुरनका समाज के वापस लौटावे के चाहत रखेला, या फेर आधुनिक उत्पादन आ विनिमय के साधनन के पुरनका संपत्ति संबंधन के ढाँचा में जबर्दस्ती बान्हे के—ओह संबंधन के जेकरा एह आधुनिक साधनन द्वारा तोडल जाना तय रहे आ तोड़ दिहल गइल। दुनो हाल में ई प्रतिक्रियावादी आ काल्पनिक दुनो ह।

एकर अंतिम शब्द बाड़ें: हस्तशिल्प में कॉरपोरेट गिल्ड, आ खेती में पितृसत्तात्मक संबंध।

अंत में, जब सख्त ऐतिहासिक सच्चाइयों आपन खुद के धोखा के सभ नशा उड़ा दिहलीं, त समाजवाद के ई रूप अवसाद आ निराशा के घिनौना दौर में खतम हो गइल।


\subsection*{{\engfont C}. जर्मन या "सच्चा" समाजवाद}
फ्रांस के समाजवादी आ कम्युनिस्ट साहित्य, जे सत्ताधारी बुर्जुआ के दबाव में पैदा भइल रहे आ जे ओह सत्ता के खिलाफ लड़ाई के अभिव्यक्ति रहे, जर्मनी में ओह समय आइल जब ओहिजा के बुर्जुआ वर्ग सामंती निरंकुशता संगे आपन लड़ाई अभी शुरूए कइले रहे।

जर्मन दार्शनिक, दार्शनिक बने के चाहत रखे वाला आ लच्छेदार बात करे वाला लोग एह साहित्य पर लपक के कब्जा कर लिहलन, बाकिर ई भुला गइलन कि जब ई फ्रांसीसी रचनाएं फ्रांस से जर्मनी आइलीं, त उनकरा संगे फ्रांसीसी सामाजिक स्थितियां जर्मनी ना आइलीं! जर्मन सामाजिक स्थितियन के संपर्क में आ के एह फ्रांसीसी साहित्य आपन सारा व्यावहारिक महत्व खो दिहलस आ महज एगो विशुद्ध साहित्यिक रूप ले लेहलस। एह तरहे 18वीं सदी के जर्मन दार्शनिकों खातिर पहिलका फ्रांसीसी क्रांति के माँग केवल सामान्य "व्यावहारिक विवेक" के माँग बन के रह गइलीं, आ क्रांतिकारी फ्रांसीसी बुर्जुआ वर्ग के इच्छा के अभिव्यक्ति उनकर आँख में शुद्ध 'इच्छा' के, अइसन इच्छा के जेकरा होखे के चाहिए रहे, यानी सामान्य सच्चा मानवीय इच्छा के कानून बन गइल!

जर्मन विद्वानन के काम बस ईहे रहे कि ऊ नया फ्रांसीसी विचारन के आपन प्राचीन दार्शनिक ज़मीर संगे मेल में ले आवें, या यूँ कहें कि आपन दार्शनिक नजरिया छोड़े बिना फ्रांसीसी विचारन के अपना लें।

ई अपनाव ओही तरहे भइल जिस तरहे कवनो बिदेसी भाखा सीखल जाला—यानी अनुवाद द्वारा!

ई सब जानत बा कि कइसे पादरी लोग पुरनका प्राचीन क्लासिक पांडुलिपियन के ऊपर कैथोलिक संतन के बेवकूफी भरी जीवनियाँ लिख देत रहलन। जर्मन विद्वान लोग धर्मनिरपेक्ष फ्रांसीसी साहित्य संगे ओकर उल्टा कइलन। ऊ लोग फ्रांसीसी मूल पाठ के नीचे आपन दार्शनिक बेवकूफी लिख दिहलन। उदाहरण खातिर, रुपिया-पैसा के आर्थिक कामन के फ्रांसीसी आलोचना के नीचे ऊ लोग लिखलन "मानवता के अलगाव", आ बुर्जुआ राज्य के फ्रांसीसी आलोचना के नीचे लिखलन "सामान्य वर्ग के तख्तापलट", इत्यादि!

फ्रांसीसी ऐतिहासिक आलोचना के पीछे एह दार्शनिक जुमलेबाज़ी के घुसावे के ऊ लोग "कर्म के दर्शन", "सच्चा समाजवाद", "जर्मन समाजवाद के विज्ञान", "समाजवाद के दार्शनिक आधार" वगैरह के नाम दिहलन।

एह तरहे फ्रांसीसी समाजवादी आ कम्युनिस्ट साहित्य के पूरा तरह नपुंसक बना दिहल गइल। आ काहे से कि जर्मन विद्वान के हाथ में आ के ई साहित्य एक वर्ग के दोसर वर्ग संगे लड़ाई के बयान कइल बंद कर दिहले रहे, एह से ऊ समझे लागल कि ऊ "फ्रांसीसी एकांगीपन" से ऊपर उठ चुकल बा आ असली ज़रूरतन के जगह सत्य के ज़रूरतन के प्रतिनिधित्व करत बा; सर्वहारा के हितन के जगह मानव स्वभाव के हितन के, सामान्य इंसान के हितन के जे कवनो वर्ग के नइखे, जेकर कवनो असली वजूद नइखे, जे बस दार्शनिक कल्पना के धुंधला दुनिया में जीएला!

ई जर्मन समाजवाद, जे आपन स्कूली लइका वाले काम के अतना गंभीरता आ गांभीर्य से लेत रहे आ आपन घटिया माल के मदारी नियर ढिंढोरा पीटत रहे, धीरे-धीरे आपन पांडित्यपूर्ण मासूमियत खो दिहलस।

जर्मन, आ खास करके प्रूशियन बुर्जुआ वर्ग के सामंती अभिजात वर्ग आ निरंकुश राजशाही के खिलाफ लड़ाई—यानी उदारवादी आंदोलन—अउरी गंभीर हो गइल।

एकरा से "सच्चा" समाजवाद के ऊ बहु-प्रतीक्षित मउका मिल गइल कि ऊ राजनीतिक आंदोलन के सामने समाजवादी माँगें रख सके, उदारवाद, प्रतिनिधि सरकार, बुर्जुआ प्रतियोगिता, प्रेस के बुर्जुआ आज़ादी, बुर्जुआ कानून, बुर्जुआ स्वतंत्रता आ समानता पर पुरनका शाप बरसा सके, आ आम जनता के उपदेश दे सके कि एह बुर्जुआ आंदोलन से उनकरा पाए खातिर कुछ नइखे, खोवे खातिर सब कुछ बा! जर्मन समाजवाद एन वक्त पर ई भुला गइल कि फ्रांसीसी आलोचना (जेकर ऊ बेवकूफ गूँज रहे) आधुनिक बुर्जुआ समाज के, ओकरा अनुकूल आर्थिक स्थितियन आ राजनीतिक संविधान के वजूद के मान के चलत रहे—अइसन चीज़ें जिनकरा हासिल कइल ही जर्मनी में चलत लड़ाई के मकसद रहे!

निरंकुश सरकारन आ उनकर पादरियन, प्रोफेसरन, जमींदारन आ नौकरशाहन खातिर ई "सच्चा समाजवाद" मँडरात बुर्जुआ वर्ग के डरावे वाला एगो मनभावन बिझुका बन गइल।

ई ओह कड़वी गोलियन (चाबुक आ गोली) के बाद एगो मीठा स्वाद रहे जेकरा से ओही सरकारें ओह समय जर्मन मज़दूरन के बगावत के दबावत रहलीं!

जहाँ ई "सच्चा" समाजवाद सरकारन खातिर जर्मन बुर्जुआ संगे लड़े के हथियार बनल, ओही संगे ई सीधा रूप से एगो प्रतिक्रियावादी हित के—जर्मन दकियानूसी छोट-बुर्जुआ के हित के—प्रतिनिधित्व करत रहे। जर्मनी में 16वीं सदी के अवशेष छोट-बुर्जुआ वर्ग, जे तब से तरह-तरह के रूपन में बार-बार सामने आवत रहे, मउजूदा हालत के असली सामाजिक आधार रहे।

एह वर्ग के बचावे के मतलब रहे जर्मनी के मउजूदा हालत के बचावल। बुर्जुआ वर्ग के औद्योगिक आ राजनीतिक वर्चस्व एकरा खातिर निश्चित विनाश के खतरा ले के आवत रहे—एक ओर पूँजी के केंद्रीकरण से, दोसर ओर क्रांतिकारी सर्वहारा के उभरला से। "सच्चा" समाजवाद के लागल कि ऊ एक तीर से दुगो शिकार कर ली! ई महामारी नियर फैल गइल।

काल्पनिक जाल के चोगा, जेकरा पर अलंकार के फूल कढ़ल रहलन आ जेकरा के बीमार भावुकता के ओस में भिगोवल गइल रहे—ई अलौकिक चोगा जेकरा में जर्मन समाजवादी आपन सुखल आ बेजान "अमर सचन" के लपेटत रहलन, एह दकियानूसी जनता के बीच आपन माल बिकावे में अद्भुत मदद कइलस। आ आपन ओर से जर्मन समाजवाद छोट-बुर्जुआ दकियानूसी के लच्छेदार प्रतिनिधि के रूप में आपन भूमिका के बेसी से बेसी स्वीकार करत गइल।

ई जर्मन राष्ट्र के आदर्श राष्ट्र आ जर्मन छोट-दकियानूसी के आदर्श इंसान घोषित कइलस। एह आदर्श इंसान के हर नीच आ घटिया हरकत के ई एगो छिपल, उच्च, समाजवादी अर्थ देवे लागल जे ओकर असली चरित्र के एकदम उल्टा रहे। ई कम्युनिज्म के "क्रूर विनाशकारी" झुकाव के सीधा विरोध करे आ सभ वर्ग-संघर्षन के प्रति आपन परम आ निष्पक्ष नफरत घोषित करे के हद तक गइल। बहुत कम अपवादों के छोड़ के, जर्मनी में (1847 में) चलत सभ तथाकथित समाजवादी आ कम्युनिस्ट रचनाएं एह घिनौना आ बेजान साहित्य के दायरा में आवेलीं।


\subsection*{{\engfont 2}. रूढ़िवादी या बुर्जुआ समाजवाद}
\addcontentsline{toc}{subsection}{{\engfont 2}. रूढ़िवादी या बुर्जुआ समाजवाद}
बुर्जुआ वर्ग के एगो हिस्सा बुर्जुआ समाज के वजूद सुरक्षित रखे खातिर सामाजिक शिकायतन के दूर करे के चाहत रखेला।

एह हिस्सा में अर्थशास्त्री, परोपकारी, मानवतावादी, मज़दूर वर्ग के हालत सुधारे वाला, परोपकार के आयोजक, पशु-क्रूरता निवारण सभा के सदस्य, नशाबंदी के कट्टर समर्थक, आ हर तरह के छुपल-छोपल सुधारक शामिल बाड़ें। समाजवाद के एह रूप के पूरा-पूरा प्रणाली में विकसित कइल गइल बा।

हमनी एह रूप के उदाहरण के तौर पर प्रूडॉन के किताब 'गरीबी के दर्शन' के नाम ले सकत बानी।

समाजवादी बुर्जुआ आधुनिक सामाजिक स्थितियन के सभ फायदा चाहेलान, बाकिर ओकरा से लाज़मी रूप से पैदा होखे वाला संघर्षन आ खतरन के बिना! ऊ मउजूदा समाज के ओकर क्रांतिकारी आ विघटनकारी तत्वन के बिना चाहेलान। ऊ सर्वहारा के बिना बुर्जुआ वर्ग चाहेलान! बुर्जुआ वर्ग स्वाभाविक रूप से ओह दुनिया के सबसे बढ़िया दुनिया समझेला जहाँ ओकर हुकूमत बा; आ बुर्जुआ समाजवाद एह आरामदेह सोच के अलग-अलग बेसी या कम पूरा प्रणालियन में विकसित करेला। सर्वहारा से अइसन प्रणाली लागू करे के आ एह तरहे सीधा सामाजिक "नया येरूशलम" में मार्च करे के माँग करत, ई असल में बस ईहे माँग करेला कि सर्वहारा मउजूदा समाज के सीमा में रहे, बाकिर बुर्जुआ वर्ग के बारे में आपन सभ नफरत भरे विचारन के उखाड़ के फेंक दे!

एह समाजवाद के दोसर आ बेसी व्यावहारिक, बाकिर कम व्यवस्थित रूप मज़दूर वर्ग के नजर में हर क्रांतिकारी आंदोलन के नीचा दिखावे के कोशिश कइलस—ई देखा के कि कवनो महज राजनीतिक सुधार ना, बलुक केवल भौतिक जीवन-स्थितियन में, यानी आर्थिक संबंधन में बदलाव ही उनकरा खातिर फायदेमंद हो सकत बा। बाकिर भौतिक जीवन-स्थितियन में बदलाव से एह समाजवाद के मतलब बुर्जुआ उत्पादन संबंधन के खतमी (जे केवल क्रांति द्वारा हो सकत बा) एकदम नइखे, बलुक एह संबंधन के बनावे रखत प्रशासनिक सुधार ह—अइसन सुधार जे पूँजी आ श्रम के बीच के संबंधन पर कवनो असर ना डालें, बलुक बेसी से बेसी बुर्जुआ सरकार के खरचा घटावें आ ओकर प्रशासनिक काम के सादा बनावें।

बुर्जुआ समाजवाद आपन सही अभिव्यक्ति तब आ केवल तब हासिल करेला जब ई महज एगो मुहावरा या अलंकार बन जाला।

मुफ़्त व्यापार: मज़दूर वर्ग के फायदा खातिर! संरक्षण शुल्क: मज़दूर वर्ग के फायदा खातिर! जेल सुधार: मज़दूर वर्ग के फायदा खातिर! ई बुर्जुआ समाजवाद के अंतिम आ एकमात्र संजीदा शब्द ह।

ओकरा एह वाक्य में समेटल जा सकत बा: बुर्जुआ एगो बुर्जुआ ह—मज़दूर वर्ग के फायदा खातिर!


\subsection*{{\engfont 3}. आलोचनात्मक-काल्पनिक समाजवाद आ कम्युनिज्म}
\addcontentsline{toc}{subsection}{{\engfont 3}. आलोचनात्मक-काल्पनिक समाजवाद आ कम्युनिज्म}
हमनी इहँवा ओह साहित्य के बात नइखी करत जे हर बड़ आधुनिक क्रांति में हमेशा सर्वहारा के माँगन के आवाज़ दिहले बा, जइसे बाबेफ आ दोसर लोगन के रचनाएं।

सर्वहारा द्वारा आपन मकसद हासिल करे के पहिलका सीधा कोशिश, जे सार्वभौमिक उत्तेजना के समय सामंती समाज के गिरावे के दौर में कइल गइल रहे, लाज़मी रूप से नाकाम भइल—काहे से कि ओह समय सर्वहारा खुद अविकसित हालत में रहे आ ओकरा आज़ादी खातिर जरूरी आर्थिक स्थितियां गायब रहलीं, अइसन स्थितियां जिनकरा पैदा होना बाकी रहे आ जे केवल आवे वाला बुर्जुआ युग द्वारा ही पैदा कइल जा सकत रहलीं। सर्वहारा के एह पहिलका आंदोलनों संगे आइल क्रांतिकारी साहित्य के चरित्र लाज़मी रूप से प्रतिक्रियावादी रहे। ई सार्वभौमिक वैराग्य आ सबसे भद्दे रूप में सामाजिक समता के उपदेश देत रहे।

असली मायने में कहल जाए वाला समाजवादी आ कम्युनिस्ट प्रणालियां—सेंट-सिमोन, फूरिए, ओवेन आ दोसर लोगन के—सर्वहारा आ बुर्जुआ के बीच के संघर्ष के ओह शुरुआती अविकसित दौर में पैदा भइलीं जेकर जिक्र हमनी ऊपर कइले बानी (देखीं भाग 1: बुर्जुआ आ सर्वहारा)।

एह प्रणालियन के संस्थापक मउजूदा समाज में वर्ग-विरोधन के संगे-संगे विघटनकारी तत्वन के असर भी देखेलान। बाकिर सर्वहारा, जे अभी आपन बचपन में बा, उनकरा सामने अइसन वर्ग के रूप में आवेला जेकरा पास कवनो ऐतिहासिक पहल या कवनो आज़ाद राजनीतिक आंदोलन नइखे।

काहे से कि वर्ग-विरोध के विकास उद्योग के विकास संगे ताल मिला के चलेला, एह से जे आर्थिक स्थिति ऊ पावेलान ऊ उनकरा सर्वहारा के आज़ादी खातिर भौतिक स्थितियां मुहैया ना करा पावे। एह से ऊ नया सामाजिक विज्ञान, नया सामाजिक कानूनन के खोज में जालें जे एह स्थितियन के पैदा कर सकें।

ऐतिहासिक कार्रवाई के उनकर व्यक्तिगत रचनात्मक कार्रवाई के सामने घुटना टेके के पड़ेला, आज़ादी के ऐतिहासिक रूप से उपजल स्थितियन के काल्पनिक स्थितियन के सामने, आ सर्वहारा के धीरे-धीरे स्वतःस्फूर्त वर्ग-संगठन के एह आविष्कारकों द्वारा खास तौर पर गढ़ल समाज के संगठन के सामने! आवे वाला इतिहास उनकर नजर में उनकर सामाजिक योजनाओं के प्रचार आ व्यावहारिक अमल में सिमट जाला।

आपन योजनाएं बनावला में ऊ एह बात के प्रति सचेत बाड़ें कि ऊ सबसे बेसी सतावल आ दुखी वर्ग के रूप में मज़दूर वर्ग के हितन के खयाल रखत बाड़ें। केवल सबसे बेसी दुखी वर्ग के नजरिया से ही सर्वहारा उनकर खातिर वजूद रखेला।

वर्ग-संघर्ष के अविकसित रूप आ उनकर आपन माहौल एह तरह के समाजवादियन के खुद के सभन वर्ग-विरोधन से बहुत ऊपर समझे पर मजबूर करेला। ऊ समाज के हर सदस्य के, इहँवा तक कि सबसे बेसी सुख-सुविधा वाले सदस्य के भी हालत सुधारे के चाहत रखेलान। एह से ऊ बिना कवनो वर्ग-भेद के पूरे समाज से, आ असल में सत्तारूढ़ वर्ग से तरजीह दे के गुहार लगावेलान। काहे से कि एक बेर लोग उनकर प्रणाली समझ लेहें, त ऊ ओकरा में सबसे बढ़िया समाज के सबसे बढ़िया योजना देखे से कइसे चूक सकत बाड़ें?

एह से ऊ सभन राजनीतिक आ खास करके सभन क्रांतिकारी कार्रवाई के खारिज कर देवेलान; ऊ आपन मकसद शांतिपूर्ण साधनन से हासिल करे के चाहत रखेलान, आ छोट-छोट प्रयोगों द्वारा (जेकर नाकाम होना तय बा) आ उदाहरण के ताक़त से नया सामाजिक सुसमाचार खातिर रास्ता साफ करे के कोशिश करेलान।

भविष्य के समाज के अइसन काल्पनिक चित्र, जे ओह समय खींचल गइल जब सर्वहारा अभी बहुत अविकसित हालत में रहे आ आपन स्थिति के केवल काल्पनिक विचार रखत रहे, समाज के पुनर्गठन खातिर ओह वर्ग के पहिलका सहज इच्छाओं संगे मेल खाएला।

बाकिर एह समाजवादी आ कम्युनिस्ट रचनाओं में एगो आलोचनात्मक तत्व भी बा। ऊ मउजूदा समाज के हर सिद्धांत पर हमला करेलान। एह से ऊ मज़दूर वर्ग के चेतना खातिर बेहद कीमती सामग्री से भरल बाड़ें। ओकरा में प्रस्तावित व्यावहारिक उपाय—जइसे गाँव आ शहर के अंतर के खतमी, परिवार के खतमी, निजी व्यक्तिन के खरचा पर उद्योग चलावे के खतमी, मजदूरी प्रणाली के खतमी, सामाजिक तालमेल के घोषणा, राज्य के काम-काज के महज उत्पादन के देख-रेख में बदलना—ई सभ प्रस्ताव केवल वर्ग-विरोधन के खतमी के ओर इशारा करेलान जे ओह समय अभी अंकुरित होत रहलन आ जिनकरा एह रचनाओं में केवल उनकर शुरुआती आ अस्पष्ट रूपन में पहचान कइल गइल रहे। एह से एह प्रस्तावों के चरित्र पूरा तरह काल्पनिक बा।

आलोचनात्मक-काल्पनिक समाजवाद आ कम्युनिज्म के महत्व के ऐतिहासिक विकास संगे उल्टा नाता होला। जइसे-जइसे आधुनिक वर्ग-संघर्ष विकसित होला आ निश्चित रूप लेवेला, ई लड़ाई से काल्पनिक रूप से अलग रहे के सोच आ ओकरा पर काल्पनिक हमला आपन सारा व्यावहारिक मूल्य आ सैद्धांतिक तर्क खो देवेलान। एह से, भले एह प्रणालियन के संस्थापक कई मायने में क्रांतिकारी रहलन, उनकर चेला लोग हर मामले में महज प्रतिक्रियावादी पंथ बन गइलन। ऊ मज़दूर वर्ग के प्रगतिशील ऐतिहासिक विकास के विरोध में आपन गुरुओं के पुरनका विचारन के मजबूती से पकड़े बाड़ें। एह से ऊ वर्ग-संघर्ष के मंद करे आ वर्ग-विरोधन के मेल-मिलाप करावे के निष्ठापूर्वक कोशिश करेलान। ऊ अभी भी आपन सामाजिक यूटोपिया के प्रायोगिक अमल के सपना देखेलान—अलग-थलग फ़ैलेन्स्टेरे बनावे के, होम कॉलोनी स्थापित करे के, छोट इकारिया खड़ा करे के—आ एह सभ हवा-हवाई महलन के बनाए खातिर ऊ बुर्जुआ वर्ग के भावना आ तिजोरी से गुहार लगावे पर मजबूर बाड़ें। धीरे-धीरे ऊ ऊपर बतलावल प्रतिक्रियावादी रूढ़िवादी समाजवादियन के श्रेणी में गिर जालें, जिनकरा से ऊ केवल बेसी व्यवस्थित पांडित्य आ आपन सामाजिक विज्ञान के जादुई असर में आपन कट्टर व अंधविश्वास भरे भरोसे से अलग बाड़ें।

एह से ऊ मज़दूर वर्ग के कवनो भी राजनीतिक कार्रवाई के उग्र विरोध करेलान; अइसन कार्रवाई उनकर नजर में केवल नया सुसमाचार में अंधे अविश्वास के नतीजा होला।

इंग्लैंड में ओवेनवादी चार्टिस्टों के विरोध करेलान, आ फ्रांस में फूरिएवादी रिफॉर्मिस्टों के।

\cleardoublepage
\section*{भाग {\engfont IV}. मउजूदा अलग-अलग विपक्षी पार्टियन संगे कम्युनिस्टों के नाता}
\addcontentsline{toc}{section}{भाग {\engfont IV}. मउजूदा अलग-अलग विपक्षी पार्टियन संगे कम्युनिस्टों के नाता}
भाग 2 में मउजूदा मज़दूर पार्टियन संगे कम्युनिस्टों के नाता साफ़ कर दिहल गइल बा, जइसे इंग्लैंड में चार्टिस्ट आ अमेरिका में कृषि सुधारक।

कम्युनिस्ट मज़दूर वर्ग के तात्कालिक मकसदन के हासिल करे, उनकर क्षणिक हितन के लागू करावे खातिर लड़ेलान; बाकिर मउजूदा आंदोलन में ऊ ओकरा भविष्य के भी प्रतिनिधित्व करेलान आ ओकर खयाल रखेलान। फ्रांस में कम्युनिस्ट लोग रूढ़िवादी आ उग्रवादी बुर्जुआ के खिलाफ सोशल-डेमोक्रेटों संगे हाथ मिलावेलान, बाकिर महान क्रांति से विरासत में आइल जुमलेबाज़ी आ भ्रमों के प्रति आलोचनात्मक नजरिया रखे के आपन अधिकार सुरक्षित रखेलान।

स्वीट्जरलैंड में ऊ रेडिकल लोग के समर्थन करेलान, बाकिर एह बात के नजरअंदाज कइले बिना कि एह पार्टी में विरोधी तत्व बाड़ें—आधा फ्रांसीसी मायने में लोकतांत्रिक समाजवादी, आ आधा उग्रवादी बुर्जुआ।

पोलैंड में ऊ ओह पार्टी के समर्थन करेलान जे राष्ट्रीय मुक्ति के पहिलका शर्त के रूप में कृषि क्रांति पर जोर देवेले—ऊ पार्टी जे 1846 के क्राको विद्रोह सुलगवले रहे।

जर्मनी में जब-जब बुर्जुआ वर्ग निरंकुश राजशाही, सामंती जमींदारन आ छोट-बुर्जुआ के खिलाफ क्रांतिकारी रूप से काम करेला, कम्युनिस्ट ओकरा संगे मिल के लड़ेलान।

बाकिर ऊ कवनो क्षण मज़दूर वर्ग के भीतर बुर्जुआ आ सर्वहारा के बीच के शत्रुतापूर्ण विरोध के सबसे साफ़ चेतना भरे के बंद ना करेलान, ताकि जर्मन मज़दूर बुर्जुआ वर्ग के खिलाफ उन सामाजिक आ राजनीतिक स्थितियन के हथियार के रूप में इस्तेमाल कर सकें जे बुर्जुआ के आपन हुकूमत संगे लावे के पड़ेला, आ ताकि जर्मनी में प्रतिक्रियावादी वर्गन के पतन के तुरंत बाद खुद बुर्जुआ वर्ग के खिलाफ लड़ाई शुरू कइल जा सके।

कम्युनिस्ट आपन मुख्य ध्यान जर्मनी पर केंद्रित करेलान, काहे से कि ऊ देस अइसन बुर्जुआ क्रांति के मुहाने पर बा जेकरा 17वीं सदी के इंग्लैंड आ 18वीं सदी के फ्रांस के मुकाबले यूरोपीय सभ्यता के बेसी विकसित स्थितियन में आ बेसी विकसित सर्वहारा संगे पूरा होना तय बा, आ काहे से कि जर्मनी में बुर्जुआ क्रांति ओकरा तुरंत बाद आवे वाली सर्वहारा क्रांति के केवल प्रस्तावना होई।

संक्षेप में, कम्युनिस्ट हर जगह मउजूदा सामाजिक आ राजनीतिक व्यवस्था के खिलाफ हर क्रांतिकारी आंदोलन के समर्थन करेलान।

एह सभन आंदोलनों में ऊ संपत्ति के सवाल के, चाहे ओह समय ओकर विकास के स्तर कवनो होखे, आंदोलन के मुख्य आ केंद्रीय सवाल के रूप में सामने ले आवेलान।

अंत में, ऊ हर जगह सभन देसन के लोकतांत्रिक पार्टियन के एकता आ सहमति खातिर काम करेलान।

कम्युनिस्ट आपन विचार आ मकसद छुपावे से नफरत करेलान। ऊ खुलेआम घोषित करेलान कि उनकर मकसद केवल सभन मउजूदा सामाजिक स्थितियन के बलपूर्वक उखाड़ फेंके से ही हासिल हो सकत बा। सत्तारूढ़ वर्गन के कम्युनिस्ट क्रांति से काँपे दीं! सर्वहारा के पास आपन जंजीरन के सिवाय खोवे खातिर कुछ नइखे। उनकरा पास जीते खातिर पूरा दुनिया बा!

दुनिया के सभे मज़दूर, एक हो जाईं!

\cleardoublepage
\section*{परिशिष्ट: प्रमुख शब्दकोश}
\addcontentsline{toc}{section}{परिशिष्ट: प्रमुख शब्दकोश}
एह घोषणापत्र में आवे वाला मुख्य राजनीतिक आ दार्शनिक शब्दन के भोजपुरी अर्थ आ व्याख्या, ताकि पाठक एह विचारन के गहराई से समझ सकें:

१. बुर्जुआ {\engfont (Bourgeois / Bourgeoisie)}

मुख्य अर्थ: पूँजीपति मालिक, कारखाना आ संपत्ति के स्वामी।

सहज अर्थ: जेकरा पास उत्पादन के मुख्य साधन बाड़ें—जइसे बड़का कारखाना, मशीन, व्यापार आ पूँजी। ई वर्ग खुद शारीरिक मेहनत ना करे, बलुक दोसरा मज़दूरन के मेहनत पर मुनाफ़ा कमाएला आ समाज पर आपन दबदबा बनाए रखेला।

२. सर्वहारा {\engfont (Proletarian / Proletariat)}

मुख्य अर्थ: मज़दूर वर्ग, श्रमजीवी समाज।

सहज अर्थ: अइसन लोग जेकरा पास आपन कवनो खेत, कारखाना या पूँजी नइखे। आपन आ आपन बाल-बच्चन के पेट पाले खातिर जेकरा रोज आपन मेहनत (श्रम) पूँजीपति मालिकन के पास बेचे के पड़ेला।

३. कम्युनिज्म / साम्यवाद {\engfont (Communism)}

मुख्य अर्थ: सबकरा खातिर बराबरी, साझी संपत्ति के समाज।

सहज अर्थ: अइसन सामाजिक व्यवस्था जहाँ धन, जमीन आ कारखानन पर कवनो गिने-चुने अमीर मालिक के कब्ज़ा ना हो के पूरा समाज आ मज़दूरन के साझी अधिकार होखे। जहाँ कवनो शोषक-शोषित ना होखे।

४. पूँजी आ पूँजीवाद {\engfont (Capital \& Capitalism)}

मुख्य अर्थ: मुनाफा कमाए वाली धन-दौलत आ ओकर व्यवस्था।

सहज अर्थ: अइसन व्यवस्था जहाँ हर चीज़ (अनाज, कपड़ा, मकान, दवाई) केवल मुनाफा कमाए आ खरीद-बिक्री खातिर बनेला। एह में मज़दूरन के खून-पसीना से उपजल मुनाफा मालिकन के तिजोरी में जाला।

५. सामंतवाद {\engfont (Feudalism)}

मुख्य अर्थ: जमींदारी व्यवस्था, राजा-महाराजा आ नवाबन के जमाना।

सहज अर्थ: पूँजीवाद से पहिले के समय जहाँ गाँव के बड़-बड़ जमींदार आ राजा लोग किसानन आ रैयत से जबरन लगान आ बेगारी असूल के शासन करत रहलन।

६. वर्ग-संघर्ष {\engfont (Class Struggle)}

मुख्य अर्थ: शोषक (दबावे वाला) आ शोषित (दबाये वाला) के बीच के लड़ाई।

सहज अर्थ: इतिहास के हर दौर में समाज में दू गो मुख्य वर्ग रहल बाड़ें—एगो मालिक आ दोसर मज़दूर। एह दुनू के हित एक-दूसरा के उल्टा होला, आ एह दुनू के बीच के खिंचाव आ लड़ाई ही वर्ग-संघर्ष ह।

७. उत्पादन के साधन {\engfont (Means of Production)}

मुख्य अर्थ: सामान पैदा करे वाला ज़रिया (जमीन, मशीन, कारखाना, खदान)।

सहज अर्थ: कवनो सामान या उपज पैदा करे खातिर लागल औजार, बड़का मशीन, जमीन, कारखाना आ कच्चा माल, जेकरा बिना कवनो निर्माण संभव नइखे।

८. क्रांति {\engfont (Revolution)}

मुख्य अर्थ: शोषक ढाँचा के उखाड़ के नया न्यायपूर्ण व्यवस्था लावल।

सहज अर्थ: समाज के सड़ल-गलल आ शोषक व्यवस्था के आम मज़दूरन द्वारा एक संगे मिल के जड़ से उखाड़ फेंकल आ ओकरा जगह सबकरा हक वाली नई व्यवस्था कायम कइल।


\end{document}
